%% Translated from IEEE conference format (manuscript-ieee/main.tex) to the
%% Springer Nature sn-jnl journal template.
%%
%% NOTE: sn-jnl.cls ships several reference-style options (sn-basic, sn-nature,
%% sn-vancouver-num, sn-mathphys-num, sn-apa, sn-chicago, ...). sn-vancouver-num
%% is used below since Vancouver numbered citation style is the medical-journal
%% norm and matches the numeric \cite{} style already used throughout this
%% document -- confirm this is the reference style required by the specific
%% journal you are submitting to and swap the option below if not.
\documentclass[pdflatex,sn-vancouver-num]{sn-jnl}

\usepackage{graphicx}
\usepackage{multirow}
\usepackage{amsmath,amssymb,amsfonts}
\usepackage{xcolor}
\usepackage{textcomp}
\usepackage{booktabs}
\let\cline\cmidrule % sn-jnl.cls aliases \cline to \cmidrule before booktabs is loaded, leaving it undefined; redo the alias here

% custom packages carried over from the IEEE version
\usepackage[font=footnotesize]{caption}
\usepackage{subcaption}
\usepackage{rotating}
\usepackage{cleveref}
\usepackage{float}
\usepackage{lineno}
\usepackage{newunicodechar}
\newunicodechar{−}{\ensuremath{-}} % Unicode MINUS SIGN (U+2212), unmapped in the default font encoding

\raggedbottom

\begin{document}

\title[ML and Counterfactuals for DPN Pre-screening]{An Application of Machine Learning and Counterfactuals for the Pre-screening and Explanation of Diabetic Peripheral Neuropathy}

\author*[1,2]{\fnm{Antonio C.} \sur{Briza}}\email{antonio.briza@asti.dost.gov.ph, AntonioBriza.PhDinDS2024@aim.edu}

\author[3]{\fnm{Niven} \sur{Fernan}}\email{email address or ORCID}
\equalcont{These authors contributed equally to this work.}

\author[3]{\fnm{Elizer M.} \sur{Montemayor}}\email{email address or ORCID}

\author[3]{\fnm{Levin} \sur{Danganan}}\email{email address or ORCID}

\author[3]{\fnm{Dan Neftalie A.} \sur{Juangco}}\email{email address or ORCID}

\author[3]{\fnm{Evelyn O.} \sur{Chua-Ley}}\email{email address or ORCID}

\author[4]{\fnm{Roy Raoul} \sur{Felipe}}\email{email address or ORCID}

\author[3]{\fnm{Deborah A.} \sur{Bernardo}}\email{email address or ORCID}

\author[3]{\fnm{Levi S.} \sur{Rejuso Jr}}\email{email address or ORCID}

\author[3]{\fnm{Marissa T.} \sur{Ong}}\email{email address or ORCID}

\author[2]{\fnm{Christopher P.} \sur{Monterola}}\email{CMonterola@aim.edu}

\affil*[1]{\orgdiv{Computer Software Division}, \orgname{DOST-Advanced Science and Technology Institute}, \orgaddress{\city{Quezon City}, \country{Philippines}}}

\affil[2]{\orgdiv{Analytics, Computing, and Complex Systems Lab}, \orgname{Asian Institute of Management}, \orgaddress{\city{Makati City}, \country{Philippines}}}

\affil[3]{\orgdiv{Department of Neurosciences}, \orgname{East Avenue Medical Center}, \orgaddress{\city{Quezon City}, \country{Philippines}}}

\affil[4]{\orgdiv{Section of Endocrinology}, \orgname{East Avenue Medical Center}, \orgaddress{\city{Quezon City}, \country{Philippines}}}

\abstract{\textbf{Background and Objectives:} Diabetic peripheral neuropathy (DPN) is the most prevalent polyneuropathy among Filipinos with diabetes, yet remains underdiagnosed in up to 81\% of affected patients. Definitive diagnosis through nerve conduction study (NCS) is costly and geographically inaccessible in many settings. This study investigates the use of machine learning (ML) and counterfactual explanation techniques as a pre-diagnostic screening tool to identify patients with a high probability of DPN and prioritize them for NCS referral.

\textbf{Methods:} A total of 190 adult Filipino patients diagnosed with Diabetes Mellitus were recruited from a tertiary hospital in Quezon City, Philippines, between 2022 and 2023. Patients underwent Michigan Neuropathy Screening Instrument (MNSI) assessment, standard neurological examination, Sudoscan sudomotor testing, and lower extremity nerve conduction study. Patients were classified as Confirmed DPN (n = 131) or Unconfirmed DPN (n = 59) based on the 2009 Toronto consensus criteria. Fourteen classification algorithms were evaluated for baseline performance using default parameters using repeated stratified k-fold cross-validation (10 repeats × 4 folds) across twelve feature subsets. The best-performing model underwent Bayesian hyperparameter optimization. Post-hoc interpretability was provided through CatBoost feature importance, SHAP analysis, and counterfactual explanations generated using the DiCE (Diverse Counterfactual Explanations) framework.

\textbf{Results:} CatBoost achieved the highest initial Area Under the Precision-Recall Curve (AUPRC) among 14 algorithms. After optimization across repeated cross-validation, it achieved an AUPRC of 0.885 (95\% CI: 0.868--0.902), representing a 29\% relative improvement over the no-skill baseline, while obtaining a mean sensitivity of 0.924 (95\% CI: 0.904--0.944). Feature set ablation revealed that neurological examination features produced the greatest performance degradation upon removal, and that the combination of Sudoscan and neurological examination findings maintained performance close to the full feature set. CatBoost feature importance and SHAP analyses converged consistently across all four cross-validation folds on decreased ankle reflex, decreased vibration sensation, Sudoscan-derived sudomotor parameters, and diabetes duration as the dominant predictors. However, while these features ranked highest in both importance and SHAP analyses, they represent measurements of neuropathic state rather than modifiable intervention targets, limiting their direct clinical actionability. Counterfactual analysis bridged this gap by identifying the minimal changes to modifiable patient features that would alter the model's predicted DPN classification: HbA1c modification emerged as the dominant actionable pathway in 91 of 97 generated counterfactuals (93.8\%), with dyslipidemia management and chronic kidney disease status identified as secondary intervention targets, appearing in 19.6\% and 23.7\% of counterfactuals respectively.

\textbf{Conclusions:} A CatBoost-based ML model trained on non-NCS clinical features demonstrated high sensitivity for DPN pre-diagnostic screening in a Philippine clinical population, with performance competitive with published benchmarks. The multi-lens analytical approach --- spanning feature ablation, feature importance, SHAP, and counterfactual analysis --- revealed that while neurological examination and Sudoscan features are the strongest predictors of DPN risk, counterfactual explanations are necessary to translate model behavior into clinically actionable guidance. These findings support the potential utility of ML-assisted DPN screening in resource-constrained settings where NCS is unavailable or inaccessible, while identifying glycemic optimization, dyslipidemia management, and nephroprotective interventions as the most actionable modifiable risk pathways. External validation in multi-center, geographically diverse cohorts is necessary before clinical deployment.}

\keywords{diabetic peripheral neuropathy, machine learning, CatBoost, pre-diagnostic screening, counterfactual explanations, explainable artificial intelligence, Sudoscan, nerve conduction study, SHAP, Philippines}

\maketitle

\section{Introduction}
%%%%%%%%%%%%%%%%%%%%%%%
Diabetic peripheral neuropathy (DPN) is the most common type of polyneuropathy affecting 42-58\% of Filipinos \cite{ropper2019adams, araneta2006hyperglycemia, dagang2016validation}.
Characterized by progressive nerve damage, DPN often presents with symptoms such as pain, numbness, and loss of sensation in the extremities, particularly the feet. Beyond its direct impact on quality of life, DPN is a major risk factor for foot ulcers, infections, and lower-limb amputations, contributing significantly to healthcare burden.

Despite being associated with reduced quality of life and debilitating complications,
DPN is underdiagnosed in up to 81\% of patients \cite{vinik2012neuropathy, ziegler2018painful}. Hurdles in the early diagnosis of DPN includes the insidious progression, inconsistent disease presentation and the need for complicated equipment for definitive diagnosis \cite{vinik2012neuropathy, ziegler2018painful, casellini2013sudoscan}.

Skin biopsy with quantitation of intraepidermal nerve fiber density is still the gold standard in detecting small fiber neuropathy, while sudorimetry is a rapidly evolving technology that can easily be integrated into clinical practice. The development of such technologies like the Sudoscan allows noninvasive and easy-to-perform tests for detection of small fiber neuropathy based on reverse iontophoresis and chronoamperometry to measure sudomotor function \cite{mayaudon:2010newsimple, vinik2015newage}.

Early detection and management of DPN are therefore critical, yet remain challenging due to the multifactorial nature of the disease and the heterogeneity of patient profiles. In recent years, machine learning (ML) has emerged as a powerful tool for modeling complex, nonlinear relationships in medical data. By leveraging large clinical datasets, ML algorithms can identify patterns and interactions among risk factors that may not be readily apparent through traditional statistical approaches. 

Numerous studies have developed ML models with good discriminatory performance (AUCs ranging from 0.64 to 0.93), employing diverse algorithms such as random forest, XGBoost, and gradient boosting. Models often incorporate clinical, biochemical, and even traditional Chinese medicine features, enhancing predictive accuracy and relevance \cite{ma2025machine,jiang2024multifeature,cui2023study,lian2023Study,Wei2024construction}. Internal validation methods like cross-validation are widely used, and some studies report calibration and decision curve analyses, supporting clinical utility \cite{cui2023study}.

Despite the  predictive performance, many ML models—particularly ensemble and deep learning methods—are often criticized as “black boxes,” limiting their adoption in clinical settings where interpretability and trust are essential. To address this limitation, the field of explainable artificial intelligence (XAI) has gained increasing attention. XAI techniques aim to make model predictions more transparent by providing insights into how input features influence outcomes. Methods such as feature importance scores, SHAP (Shapley Additive Explanations), and partial dependence plots have been widely used to offer global and local interpretability. 

The application of XAI techniques such as SHAP, LIME, and Anchor explanations is increasingly common, providing transparent insights into feature importance and model decisions, which is critical for clinical trust and adoption \cite{wu2025prognostic,jiang2024multifeature, pandey2024decoding, wu2024interpretable, sheikh2025predicting}. Integration of XAI with multimodal data enables identification of key risk factors and supports personalized treatment recommendations \cite{nguyen2024interpretable, alHudaibi2025explainable}. Some frameworks incorporate human-in-the-loop feedback to enhance interpretability and clinical relevance \cite{lalwani2025generative}. While these approaches help elucidate model behavior, they primarily describe associations rather than actionable insights, which are crucial for clinical decision-making.

Counterfactual explanations represent a more recent and clinically meaningful approach within XAI. Instead of merely explaining why a prediction was made, counterfactuals answer “what-if” questions by identifying minimal changes to a patient’s features that would lead to a different predicted outcome. For example, a counterfactual explanation might suggest how modifying certain clinical variables—such as glycemic control or body mass index—could reduce a patient’s predicted risk of developing DPN. This aligns closely with the needs of personalized medicine, where interventions are tailored to individual patients based on modifiable risk factors. Several methodological approaches for generating counterfactual explanations have been proposed in the literature. One of the most popular is the DiCE framework \cite{mothilal2020explaining}, an optimization-based method, which formulate counterfactual generation as a constrained optimization problem that seeks minimal changes to input features while achieving a desired prediction;

The integration of machine learning and counterfactual explanations offers a promising framework for both the prediction and interpretation of diabetic peripheral neuropathy. By combining high predictive accuracy with actionable, patient-specific insights, this approach has the potential to support clinicians in early diagnosis, risk assessment, and treatment planning. Moreover, it bridges the gap between advanced computational models and practical clinical utility, fostering greater trust and adoption of AI-driven tools in healthcare. 

This study aims to investigate the application of machine learning models in conjunction with counterfactual explanation techniques for the prediction and interpretation of diabetic peripheral neuropathy (DPN) risk, thereby contributing to the growing body of research on interpretable and actionable artificial intelligence in medicine. The study also examines the potential insights derived from incorporating nerve conduction study results and Sudoscan measurements—modalities that are comparatively underrepresented in many AI-based DPN studies, which often rely primarily on routine clinical and demographic variables. Furthermore, by utilizing data collected within the Philippine clinical context, this work enhances the geographic diversity of current studies and provides context-specific insights into the application of explainable AI for DPN.


\section{Data Collection}
%%%%%%%%%%%%%%%%%%%%%%%%%%%%%
A total of 201 adult Filipino patients diagnosed with Diabetes Mellitus (DM) were recruited
between 2022 and 2023. Subjects with limb amputation, gangrenous/ necrotic foot, with implanted
cardiac pacemakers or defibrillators, and patients with other known causes of polyneuropathy were
excluded. Remaining 190 qualified subjects underwent a series of detailed assessment comprising of
Michigan Neuropathy Screening Instrument (MNSI) test, serial neurologic tests, and lower extremity NCV
examination to detect the presence of DPN. Subjects were then classified based on the 2009 Toronto
consensus criteria as confirmed, probable, possible or no DPN \cite{tesfaye2010diabetic}. This study was approved by the local Institutional Ethics Review Board and all subjects gave their written consent prior to participating.

Neuropathic symptoms were documented by subjectively asking patients for presence of any
distal extremity numbness, paresthesia, or pain, and by asking them to answer the Filipino validated
version of MNSI \cite{dagang2016validation}. Abnormal neurologic signs were detected using the standard 10g Semmes Weinstein monofilament examination (SWME) and standard 128-Hz clanging tuning fork tests \cite{feng2009semmes,oyer2007quantitative}. Additional examination for pinprick sensation was done using the same criteria for SWME. Achilles tendon reflex testing was also done and dichotomized subjects as either normal or hypo/areflexive. All neurological examinations were conducted by a single trained physician.

Nerve conduction studies were performed at a stable skin temperature of 32°C and a room
temperature of 24°C using Nihon Kohen EMG equipment with a filter width of 20 Hz to 3 kHz, amplifier
setting at 200 $\mu$V/div, and sweep speed of 10 ms/div. 

The following nerve attributes were measured: 
\begin{itemize}
\itemsep0em 
\item Sural sensory nerve action potentials (SNAP) and conduction velocities (CV)
\item Superficial peroneal SNAP and CV
\item Tibial motor nerve compound muscle action potentials (CMAP), CV, distal latencies (DL) and
F-wave latencies
\end{itemize}

All NCV testing was done by a single trained technician on the same day as other
neurological examinations. Finally, sudomotor testing using Sudoscan was performed on each participant. Subjects placed palms of their hands and the soles of their feet on stainless steel electrodes and an incremental low direct voltage (<4V) was applied for about 2 minutes. The device produced an average electrochemical skin conductance (ESC) result for both hands and feet which is based on stimulation of sweat glands by the low-level voltage, allowing evidence of sweat dysfunction \cite{garcia2022detection}.

The patients were classified into four categories: Confirmed, Probable, Possible and No DPN.  Given the limited number of samples from the latter three categories, this study combined the last three categories to form a larger group  labeled as Unconfirmed. Table~\ref{tab:data_categ} shows the breakdown among the different categories or labels. The full list of features, their groupings, description and code are shown in Table~\ref{tab:clinical_data}.

%=======================================
% Data and Categories
%=======================================
\begin{table}[h!]
\centering
\small
\begin{tabular}{lrr}
\toprule
Label & \multicolumn{2}{c}{Support}  \\
\midrule
Negative  &  9 \\
Possible  & 22 \\
Probable  & 28  \\
Unconfirmed (Negative+Possible+Probable) & &59 \\
Confirmed & & 131 \\
Overall Total & & 190 \\
\bottomrule
\end{tabular}
\caption{Distribution of patient classifications in the study dataset (N = 190). Patients were categorized based on nerve conduction study (NCS) findings into four diagnostic groups: Confirmed DPN (n = 131), Probable DPN (n = 28), Possible DPN (n = 22), and Negative DPN (n = 9). As this study aims to serve as a pre-diagnostic screening tool for DPN, the latter three groups — whose DPN status could not be definitively established without further electrodiagnostic evaluation — were consolidated into a single composite class labeled Unconfirmed (n = 59). This two-class scheme — Confirmed versus Unconfirmed — reflects the clinical objective of the model: to identify patients with a high probability of DPN and prioritize them for referral to nerve conduction studies for confirmatory diagnosis, particularly in settings where NCS may be costly or inaccessible.}
\label{tab:data_categ}
\end{table}

% =======================================
% Table of Features, Codes, Description
% =======================================
\begin{table}[h!]
\centering
\footnotesize
\begin{tabular}{llll}
\toprule
\textbf{A. Variables used in modeling} & & & \\
\midrule
\hspace{0.3cm} Group & Code & Description  & Type\\
\midrule
\hspace{0.3cm} Profile      & SEX & Patient's sex (M or F) & binary \\
             & AGE & Patient's age  & continuous \\
             & SUBJ & Other symptoms  & binary \\
             & DM DUR & Diabetes Mellitus duration (years) & continuous \\
             & INSULIN & Patient is taking insulin treatments & binary \\
             & HBA1C & Hemoglobin A1c & continuous  \\
\midrule
\hspace{0.3cm} Commorbidity & HPN & Hypertension  & binary \\
             & PAOD & Peripheral Arterial Occulsive Disease  & binary \\
             & DSLPDMIA & Dyslipidemia & binary \\
             & CKD & Chronic Kidney Disease  & binary \\
             & GBS & Guillain-Barre Syndrome  & binary \\
\midrule
\hspace{0.3cm} Neurology Test   & DEC\_VS & Decreased Vibration Sensation & binary \\
                 & DEC\_PPS & Decreased Pinprick Sensation  & binary \\
                 & DEC\_LTS & Decreased Light Touch Sensation  & binary \\
                 & DEC\_AR & Decreased Ankle Reflex  & binary \\
\midrule
\hspace{0.3cm} MNSI & MNSI & Michigan Neuropathy Screening Instrument  & continuous \\
\midrule
\hspace{0.3cm} Sudoscan & FEET\_MEAN\_ESC & Feet Mean Electrochemical Skin Conductance & continuous  \\
         & FEET\_PCT\_ASYM & Feet Percentage of Asymmetry (\%) & continuous  \\
         & HAND\_MEAN\_ESC & Hands Mean Electrochemical Skin Conductance & continuous  \\
         & HAND\_PCT\_ASYM & Hands Percentage of Asymmetry (\%) & continuous  \\
         & NS & Neuropathy Score & continuous  \\
         & CAS & Cardiac Autonomic Neuropathy Score (\%) & continuous  \\
 &  &   & \\
\midrule
\textbf{B. Variables not used in modeling} & & & \\
\midrule
\hspace{0.3cm} Group & Code & Description  & Type\\
\midrule
\hspace{0.3cm} Nerve Conduction Study & SSA\_L & Left Sural Sap Nerve Sensory Amplitude (Uv)  & continuous \\
 & SSC\_L & Left Sural Nerve Sensory Conduction (m/s)  & continuous \\
 & SPSA\_L & Left Superficial Peroneal Nerve Sensory Amplitude (Uv)  & continuous \\
 & SPSC\_L & Left Superficial Peroneal Sensory Conductivity Velocity (m/s)  & continuous \\
 & MCV\_L & Left Posterior Tibial Nerve Conduction Velocity  & continuous \\
 & DL\_L & Left Posterior Tibial Nerve Distal Latency  & continuous \\
 & CMAPANK\_L & Left Ankle Compound Muscle Action Potential  & continuous \\
 & CMAPKNE\_L & Left Knee Compound Muscle Action Potential  & continuous \\
 & FWAVE\_L & Left F-wave latency  & continuous \\
 & SSA\_R & Right Sural sap Nerve Sensory Amplitude (Uv)  & continuous \\
 & SSC\_R & Right Sural Nerve Sensory Conduction (m/s)  & continuous \\
 & SPSA\_R & Right Superficial Peroneal Nerve Sensory Amplitude (Uv) & continuous  \\
 & SPSC\_R & Right Superficial Peroneal Sensory Conductivity Velocity (m/s) & continuous  \\
 & MCV\_R & Right Posterior Tibial Nerve Conduction Velocity & continuous  \\
 & DL\_R & Right Posterior Tibial Nerve Distal Latency & continuous  \\
 & CMAPANK\_R & Right Ankle Compound Muscle Action Potential & continuous  \\
 & CMAPKNE\_R & Right Knee Compound Muscle Action Potential & continuous  \\
 & FWAVE\_R & Right F-wave Latency & continuous  \\
\bottomrule
\end{tabular}
\caption{
 Complete listing of clinical variables collected in this study, organized by feature group, with their corresponding code, description, and data type. The Profile group captures patient demographic and metabolic baseline characteristics, including sex, age, diabetes mellitus duration, insulin use, and HbA1c. The Commorbidity group encodes the presence or absence of five concurrent conditions: hypertension (HPN), peripheral arterial occlusive disease (PAOD), dyslipidemia (DSLPDMIA), chronic kidney disease (CKD), and Guillain-Barré syndrome (GBS). The Neurology Test group contains binary findings from a standard neurological examination assessing vibration, pinprick, light touch, and ankle reflex responses. The MNSI group contains the score from the Michigan Neuropathy Screening Instrument, a validated clinical tool for peripheral neuropathy screening. The Sudoscan group comprises continuous measurements of electrochemical skin conductance and asymmetry for both the hands and feet, along with a composite Neuropathy Score and a Cardiac Autonomic Neuropathy Score. The Nerve Conduction Study (NCS) group contains electrophysiological measurements of sensory amplitude, conduction velocity, distal latency, and F-wave latency across bilateral lower limb nerves; these variables served as the diagnostic reference for DPN classification but were excluded from all predictive modeling, as the study objective is to enable pre-diagnostic screening in settings where NCS is unavailable or inaccessible. Binary variables are encoded as 0 (absent) or 1 (present); continuous variables retain their original measurement units.
% A listing of all clinical variables collected in this study, including those that were used for machine learning modeling (Part A, 22 variables in total) and those that were not (Part B). Each variable is placed under a group and identified by its code, a plain-language description, and its data type (binary or continuous).  Part B of the table come from, Nerve Conduction Study — the reference standard for Diabetic Peripheral Neuropathy (DPN) diagnosis. These were excluded from machine learning modeling and are included solely for contextual reference.  The remaining features represent non-invasive or minimally invasive clinical assessments routinely available in outpatient settings, underscoring the practical deployability of any model trained on this feature space.
}
\label{tab:clinical_data}
\end{table}

Statistical tables of the data are shown for categorical and continuous features in \Cref{tab:categorical,tab:continuous} while plots are shown for categorical features, profile plus MNSI, Nerve Conduction Studies, and Sudoscan in \Cref{fig:categ,fig:profile_mnsi,fig:ncs,fig:sudo}. 


% =======================================
% Categorical Table
% =======================================
\begin{table}[h!]
\centering
\footnotesize
\begin{tabular}{lccccccccccccc}
\toprule
& \multicolumn{4}{c}{Confirmed} && \multicolumn{4}{c}{Unconfirmed} && \multicolumn{3}{c}{Statistics} \\       
\cline{2-5} \cline{7-10} \cline{12-14}
Feature     & No & Yes  & Total & Percent && No & Yes  & Total  & Percent && chi² & p(chi²)  & Cramér's V \\
\midrule
SEX & 44 & 83 & 127 & 66.8 && 16 & 47 & 63 & 33.2 && 1.27     & 0.2604 & 0.082 \\
SUBJ & 18 & 10 & 28 & 14.7 && 42 & 120 & 162 & 85.3 && 14.53  & 0.0001 & 0.277 \\
INSULIN & 42 & 69 & 111 & 58.4 && 18 & 61 & 79 & 41.6 && 4.17 & 0.0412 & 0.148 \\
\midrule
HPN         & 25 & 51   & 76    & 40.0    && 35 & 79   & 114    & 60.0    && 0.03 & 0.8734 & 0.012 \\
PAOD        & 59 & 123  & 182   & 95.8    && 1  & 7    & 8      & 4.2     && 0.64 & 0.4251 & 0.058 \\
DSLPDMIA    & 36 & 57   & 93    & 48.9    && 24 & 73   & 97     & 51.1    && 3.66 & 0.0556 & 0.139 \\
CKD         & 53 & 100  & 153   & 80.5    && 7  & 30   & 37     & 19.5    && 2.72 & 0.0991 & 0.120 \\
GBS         & 59 & 129  & 188   & 98.9    && 1  & 1    & 2      & 1.1     && 0.00 & 1.0000 & 0.000 \\
\midrule
DEC\_VS & 48 & 44 & 92 & 48.4 && 12 & 86 & 98 & 51.6 && 33.19   & 0.0000 & 0.418 \\
DEC\_PPS & 44 & 53 & 97 & 51.1 && 16 & 77 & 93 & 48.9 && 16.14  & 0.0001 & 0.291 \\
DEC\_LTS & 43 & 57 & 100 & 52.6 && 17 & 73 & 90 & 47.4 && 11.65 & 0.0006 & 0.248 \\
DEC\_AR & 51 & 45 & 96 & 50.5 && 9 & 85 & 94 & 49.5 && 39.70    & 0.0000 & 0.457 \\
\bottomrule
\end{tabular}
% \captionsetup{font=footnotesize}
\caption{Frequency distribution of categorical features stratified by binary outcome.
} 
\label{tab:categorical}
\end{table}

% =======================================
% Continuous Table
% =======================================
\begin{table}[t]
\centering
\footnotesize
\begin{tabular}{lccccccccccc}
\toprule
& \multicolumn{2}{c}{Range} &&  \multicolumn{3}{c}{Mean ± STD} && \multicolumn{4}{c}{Distribution Statistics} \\       
\cline{2-3} \cline{5-7}  \cline{9-12}
Feature & Min & Max  && Unconfirmed & Confirmed & All &&  Skewness & Kurtosis & p (MW-U) & r (effect size) \\
\midrule
AGE & 20 & 89 && 51.45 ± 13.91 & 59.82 ± 9.81 & 57.17 ± 11.89 && -0.53 & 0.44 & 0.0002 & 0.342 \\
DM\_DUR & 0 & 30 && 5.42 ± 7.51 & 8.92 ± 7.71 & 7.81 ± 7.80 && 1.24 & 0.86 & 0.0001 & 0.359 \\
HBA1C & 0 & 17.60 && 7.76 ± 2.48 & 8.39 ± 2.81 & 8.19 ± 2.72 && 0.76 & 0.59 & 0.2556 & 0.103 \\
\midrule
MNSI & 0 & 12 && 4.33 ± 2.77 & 5.76 ± 2.54 & 5.31 ± 2.69 && 0.18 & -0.59 & 0.0003 & 0.322 \\
\midrule
FEET\_MEAN\_ESC & 6.00 & 87.00 &  & 58.22 ± 16.49 & 46.32 ± 19.78 & 50.07 ± 19.56 &  & -0.35 & -0.71 & 0.0001 & -0.353 \\
FEET\_PCT\_ASYM & 0.00 & 42.00 &  & 5.52 ± 8.10 & 8.16 ± 6.94 & 7.33 ± 7.40 &  & 1.51 & 2.49 & 0.0005 & 0.315 \\
HAND\_MEAN\_ESC & 8.00 & 87.00 &  & 61.33 ± 16.94 & 54.04 ± 19.45 & 56.34 ± 18.96 &  & -0.47 & -0.59 & 0.0163 & -0.217 \\
HAND\_PCT\_ASYM & 0.00 & 54.00 &  & 9.50 ± 9.99 & 11.37 ± 9.22 & 10.78 ± 9.48 &  & 1.42 & 2.27 & 0.0767 & 0.160 \\
NS & 0.00 & 99.00 &  & 62.95 ± 17.88 & 50.98 ± 13.42 & 54.76 ± 15.94 &  & -0.01 & 1.11 & 0.0000 & -0.423 \\
CAS & 0.00 & 49.00 &  & 26.22 ± 10.87 & 31.02 ± 7.62 & 29.50 ± 9.03 &  & -0.69 & 1.00 & 0.0032 & 0.266 \\
\midrule
SSA\_L & 0.00 & 49.58 &  & 17.69 ± 6.98 & 6.45 ± 5.49 & 10.00 ± 7.95 &  & 1.09 & 2.58 & 0.0000 & -0.850 \\
SSA\_R & 0.00 & 55.41 &  & 17.88 ± 7.64 & 6.65 ± 6.99 & 10.20 ± 8.89 &  & 1.56 & 4.58 & 0.0000 & -0.838 \\
SSC\_L & 0.00 & 64.50 &  & 50.12 ± 4.79 & 33.28 ± 21.43 & 38.60 ± 19.55 &  & -1.32 & 0.08 & 0.0000 & -0.551 \\
SSC\_R & 0.00 & 59.60 &  & 49.34 ± 4.47 & 33.64 ± 21.03 & 38.60 ± 19.02 &  & -1.38 & 0.29 & 0.0000 & -0.532 \\
SPSA\_L & 0.00 & 49.79 &  & 13.66 ± 8.35 & 4.53 ± 5.68 & 7.41 ± 7.86 &  & 1.89 & 5.89 & 0.0000 & -0.720 \\
SPSA\_R & 0.00 & 46.48 &  & 13.02 ± 6.96 & 4.50 ± 4.76 & 7.19 ± 6.81 &  & 1.46 & 4.73 & 0.0000 & -0.738 \\
SPSC\_L & 0.00 & 68.50 &  & 47.93 ± 8.66 & 26.55 ± 22.73 & 33.31 ± 21.80 &  & -0.75 & -1.17 & 0.0000 & -0.592 \\
SPSC\_R & 0.00 & 65.50 &  & 47.50 ± 8.06 & 26.87 ± 22.76 & 33.39 ± 21.59 &  & -0.80 & -1.13 & 0.0000 & -0.570 \\
MCV\_L & 0.00 & 58.00 &  & 44.99 ± 4.01 & 38.73 ± 6.81 & 40.71 ± 6.72 &  & -2.26 & 12.53 & 0.0000 & -0.647 \\
MCV\_R & 0.00 & 54.00 &  & 45.92 ± 4.16 & 38.66 ± 7.18 & 40.95 ± 7.21 &  & -2.13 & 9.57 & 0.0000 & -0.704 \\
DL\_L & 0.00 & 45.50 &  & 3.80 ± 0.55 & 4.70 ± 3.76 & 4.42 ± 3.15 &  & 11.75 & 150.93 & 0.0000 & 0.449 \\
DL\_R & 0.00 & 15.15 &  & 3.82 ± 0.65 & 4.26 ± 1.43 & 4.12 ± 1.25 &  & 4.04 & 34.27 & 0.0032 & 0.266 \\
CMAPANK\_L & 0.00 & 24.80 &  & 12.25 ± 3.96 & 5.45 ± 3.42 & 7.60 ± 4.79 &  & 0.62 & 0.14 & 0.0000 & -0.822 \\
CMAPANK\_R & 0.00 & 26.74 &  & 12.11 ± 4.15 & 5.71 ± 3.37 & 7.73 ± 4.69 &  & 0.62 & 0.69 & 0.0000 & -0.787 \\
CMAPKNE\_L & 0.00 & 19.17 &  & 8.73 ± 2.67 & 3.83 ± 2.41 & 5.38 ± 3.38 &  & 0.62 & 0.54 & 0.0000 & -0.849 \\
CMAPKNE\_R & 0.00 & 19.70 &  & 8.57 ± 2.80 & 3.98 ± 2.40 & 5.43 ± 3.31 &  & 0.61 & 1.06 & 0.0000 & -0.826 \\
FWAVE\_L & 0.00 & 74.80 &  & 46.35 ± 4.19 & 51.69 ± 11.33 & 50.00 ± 9.96 &  & -2.28 & 11.44 & 0.0000 & 0.581 \\
FWAVE\_R & 0.00 & 71.10 &  & 46.50 ± 4.71 & 49.69 ± 14.17 & 48.69 ± 12.09 &  & -2.50 & 8.59 & 0.0000 & 0.512 \\
\bottomrule
\end{tabular}
% \captionsetup{font=footnotesize}
\caption{Statistics for continuous features.}
\label{tab:continuous}
\end{table}

% Categorical Plots
%=======================================
\begin{figure*}[h!]
    % \centering
    \hspace*{0.085\textwidth}
    \includegraphics[width=0.5\textwidth]{figures/eda/categ_profile_histogram.png}
    \\[2ex] % Adds a line break and some extra vertical space
    \hspace*{0.085\textwidth}
    \includegraphics[width=0.83\textwidth]{figures/eda/categ_comorbs_histogram.png}
    \\[2ex] % Adds a line break and some extra vertical space
    \hspace*{0.085\textwidth}
    \includegraphics[width=0.66\textwidth]{figures/eda/categ_neuro_histogram.png}
    \caption{
    Count distributions of categorical clinical variables comparing Unconfirmed (blue) and Confirmed (red) groups across demographic, comorbidity, and diagnostic subcategory features. Variables are organized into three rows: (1) patient profile — Sex (SEX), Subjective Neuropathy Complaint (SUBJ), and Insulin Use (INSULIN); (2) comorbidities — Hypertension (HPN), Peripheral Arterial Occlusive Disease (PAOD), Dyslipidemia (DSLPDMIA), Chronic Kidney Disease (CKD), and Guillain-Barré Syndrome (GBS); and (3) neurological test — Decreased Vibration Sense (DEC\_VS), Decreased Pinprick Sensation (DEC\_PPS), Decreased Light Touch Sensation (DEC\_LTS), and Decreased Ankle Reflex (DEC\_AR).
    Across nearly all variables, the Confirmed group consistently exhibited higher counts in the affirmative ("Yes") category relative to the Unconfirmed group, suggesting a broader and more severe clinical profile. Notably, the Confirmed group showed a strong male predominance, a substantially higher rate of subjective neuropathy complaints, and greater insulin use, potentially reflecting more advanced or insulin-dependent diabetes. Among comorbidities, HPN, DSLPDMIA, and CKD showed markedly higher affirmative counts in the Confirmed group, while PAOD and GBS affirmative cases were nearly absent in both groups, suggesting these may be rare in the study population. In the diagnostic subcategory variables, the Confirmed group consistently reported higher rates of sensory and reflex deficits across all four measures (DEC\_VS, DEC\_PPS, DEC\_LTS, DEC\_AR), consistent with a more pronounced neuropathic presentation and supporting the clinical validity of the group classification.}
    \label{fig:categ}
\end{figure*}

% Histogram of Profile MSI 
% =======================================
\begin{figure*}[h!]
    \centering
    \includegraphics[width=0.66\textwidth]{figures/eda/cont_profile_mnsi_distributions.png}
    \\[2ex] % Adds a line break and some extra vertical space
    \includegraphics[width=0.66\textwidth]{figures/eda/cont_profile_mnsi_boxplots.png}
    \caption{Stacked density distributions (top) and boxplots (bottom) of clinical and demographic variables comparing Unconfirmed (blue) and Confirmed (red) groups. Variables presented include Age (AGE; median: 52.0 vs. 61.0), Diabetes Duration in years (DM\_DUR; median: 2.0 vs. 7.0), Glycated Hemoglobin (HBA1C; median: 7.1 vs. 7.6), and Michigan Neuropathy Screening Instrument score (MNSI; median: 4.0 vs. 6.0). Three of the four variables showed statistically significant group differences (p < 0.001 for AGE, DM\_DUR, and MNSI), while HBA1C did not reach significance (p = 0.256), suggesting that glycemic control levels alone may not distinguish the two groups. The Confirmed group was notably older and had a substantially longer diabetes duration, consistent with the known association between prolonged diabetes and increased neuropathy risk. The wider interquartile range and heavier right tail observed in DM\_DUR for the Confirmed group further suggests considerable heterogeneity in disease chronicity within this group. The significantly higher MNSI scores in the Confirmed group align with expected clinical presentation, reinforcing MNSI as a potentially useful screening variable, while the marginal HBA1C difference warrants caution against over-relying on glycemic markers alone in neuropathy risk stratification.
    \label{fig:profile_mnsi}
    }
\end{figure*}

% Distribution and Boxplot of NCS 
%=======================================
\begin{figure*}[h!]
    \centering
    \includegraphics[width=1\textwidth]{figures/eda/cont_ncs1_distributions.png}
    \\[2ex] % Adds a line break and some extra vertical space
    \includegraphics[width=1\textwidth]{figures/eda/cont_ncs1_boxplots.png}
    \\[6ex] % Adds a line break and some extra vertical space
    \includegraphics[width=1\textwidth]{figures/eda/cont_ncs2_distributions.png}
    \\[2ex] % Adds a line break and some extra vertical space
    \includegraphics[width=1\textwidth]{figures/eda/cont_ncs2_boxplots.png}
    \\[6ex] % Adds a line break and some extra vertical space
    \includegraphics[width=1\textwidth]{figures/eda/cont_ncs3_distributions.png}
    \\[2ex] % Adds a line break and some extra vertical space
    \includegraphics[width=1\textwidth]{figures/eda/cont_ncs3_boxplots.png}
    \caption{
    Density distributions and boxplots of Nerve Conduction Study (NCS) variables comparing Unconfirmed (blue) and Confirmed (red) groups, encompassing Sural Peak Sensory Conduction (SPSC\_L, SPSC\_R), Motor Conduction Velocity (MCV\_L, MCV\_R), Distal Latency (DL\_L, DL\_R), Sural Sensory Amplitude (SSA\_L, SSA\_R), Sural Sensory Conduction (SSC\_L, SSC\_R), and Sural Peak Sensory Amplitude (SPSA\_L, SPSA\_R). All variables demonstrated statistically significant group differences (p < 0.001 for most; DL\_R: p = 0.003). The Confirmed group consistently showed lower median values across all conduction velocity and sensory amplitude measures, with density distributions revealing a pronounced clustering near zero in SPSC and MCV variables, indicative of a substantial subset with absent or severely reduced nerve conduction responses. Sensory amplitude variables were most dramatically reduced in the Confirmed group, with compressed interquartile ranges approaching zero across SSA and SPSA measures bilaterally, strongly suggestive of axonal loss or severe sensory nerve compromise. Modest but consistent increases in distal latency in the Confirmed group further corroborate slowed nerve conduction, collectively underscoring the discriminative utility of NCS variables in characterizing neuropathic severity between groups.
    \label{fig:ncs}
    }
\end{figure*}

% Distribution and Boxplot of Sudoscan 
%=======================================
\begin{figure*}[ht]
    \centering
    \includegraphics[width=1\textwidth]{figures/eda/cont_sudo_distributions.png}
    \\[2ex] % Adds a line break and some extra vertical space
    \includegraphics[width=1\textwidth]{figures/eda/cont_sudo_boxplots.png}
    \caption{Stacked density distributions (top) and boxplots (bottom) of all SUDOSCAN-derived variables comparing Unconfirmed (blue) and Confirmed (red) groups. Variables presented include Feet Mean Electrochemical Skin Conductance (FEET\_MEAN\_ESC; median: 60.5 vs. 47.5), Feet Percent Asymmetry (FEET\_PCT\_ASYM; median: 3.0 vs. 6.5), Hand Mean ESC (HAND\_MEAN\_ESC; median: 63.0 vs. 56.0), Hand Percent Asymmetry (HAND\_PCT\_ASYM; median: 6.0 vs. 10.0), Nerve Score (NS; median: 60.5 vs. 50.5), and Composite Autonomic Score (CAS; median: 28.0 vs. 31.0). Nearly all variables demonstrated statistically significant group differences (p < 0.001 for FEET\_MEAN\_ESC, FEET\_PCT\_ASYM, and CAS; p= 0.003 for NS; p = 0.016 for HAND\_MEAN\_ESC), with the exception of p = 0.077 for HAND\_PCT\_ASYM) potentially reflecting greater lateralized nerve impairment in the hands. Notably, the Confirmed group exhibited lower mean ESC values in both feet and hands, consistent with reduced sudomotor function. The elevated CAS in the Confirmed group, despite lower NS values, warrants further investigation into the composite scoring methodology and its constituent variables.
    \label{fig:sudo}
    }
\end{figure*}


\section{Methodology}
%%%%%%%%%%%%%%%%%%%%%%
This section describes the machine learning model development pipeline and the counterfactual generation methodology.

\subsection{Machine Learning Modeling}
In this section, we describe the machine learning models used, stratified cross-validation, feature selection, and hyperparameter optimization.

\subsubsection{Models}
%%%%%%%%%%%%%%%%%%%%%%%

A diverse set of machine learning algorithms was evaluated in this study using default or standard parameter settings to ensure a consistent and unbiased baseline comparison. These models can be broadly grouped according to their underlying learning paradigms:

\textit{(a) Linear and Probabilistic Models}.
This group includes Logistic Regression, Linear Discriminant Analysis (LDA), Quadratic Discriminant Analysis (QDA), Stochastic Gradient Descent (SGDClassifier), Naive Bayes (GaussianNB), and Linear Support Vector Machine (Linear SVM). These models are generally characterized by strong interpretability and computational efficiency. Logistic Regression and Linear SVM are effective in high-dimensional spaces and provide robust linear decision boundaries. LDA and QDA model class-conditional distributions, with QDA allowing nonlinear decision boundaries through class-specific covariance matrices. Naive Bayes is particularly efficient and performs well under conditional independence assumptions, while SGDClassifier enables scalable optimization for large datasets.

\textit{(b) Distance-Based and Instance-Based Learning}.
The k-Nearest Neighbors (kNN) algorithm falls under this category. It is a non-parametric method that makes predictions based on the proximity of data points in feature space. Its strength lies in its simplicity and ability to capture complex decision boundaries without explicit model training, although it may be sensitive to feature scaling and noise.

\textit{(c) Tree-Based and Ensemble Methods}.
This group includes Decision Tree, Random Forest, Extra Trees, and Gradient Boosting. Decision Trees are intuitive and capable of capturing nonlinear relationships and feature interactions. Random Forest and Extra Trees extend this capability through ensemble learning, improving generalization and reducing variance via bagging and randomized feature selection. Gradient Boosting builds models sequentially to correct previous errors, often achieving strong predictive performance.

\textit{(d) Advanced Gradient Boosting Frameworks}.
More sophisticated implementations of boosting are represented by XGBoost, LightGBM, and CatBoost. These algorithms are optimized for efficiency and performance on structured/tabular data. XGBoost incorporates regularization and efficient handling of sparse data; LightGBM employs histogram-based learning and leaf-wise tree growth for speed and scalability; CatBoost is particularly effective in handling categorical variables and mitigating overfitting through ordered boosting techniques.

\textit{(e) Kernel-Based Models}.
The RBF Support Vector Machine (RBF SVM) extends linear SVM by applying a radial basis function kernel, enabling the modeling of nonlinear decision boundaries. This makes it powerful for complex datasets where linear separability is not sufficient.

\textit{(f) Baseline Model}.
A Naive classifier that predicts the majority class was included as a baseline to contextualize the performance of all other models.

Overall, this diverse selection of models allows for a comprehensive evaluation across linear, nonlinear, probabilistic, and ensemble-based learning paradigms, each with distinct strengths in terms of interpretability, flexibility, robustness, and predictive power.

\subsubsection{Repeated Stratified k-Fold Cross-Validation}

Given the relatively limited number of unconfirmed cases in the dataset, a robust evaluation strategy was used to generate reliable performance estimates. A repeated stratified k-fold cross-validation scheme was employed, with 4 folds and 10 repetitions, resulting in a total of 40 validation runs per model.

Stratification ensures that each fold preserves the class distribution of the dataset, which is particularly important in imbalanced classification settings. Repetition further reduces variance in performance estimates by averaging results across multiple random partitions.



\subsubsection{Performance Metrics}
For each validation run, key performance metrics were recorded, including Receiver Operating Characteristic Area Under the Curve (ROC-AUC), Youden Index, Sensitivity (Recall), and Specificity. These metrics were selected to provide a balanced evaluation of discriminative ability and clinical relevance, particularly in distinguishing between Confirmed and Unconfirmed cases.

\subsubsection{Feature Sets}

To assess the contribution of different feature subsets, a structured feature selection and ablation strategy was implemented.

First, all models were trained using the complete set of available features to establish baseline performance. Subsequently, based on the feature groupings defined in Table~\ref{tab:clinical_data}-A, , models were retrained while systematically removing each feature group. This ablation approach enables the identification of feature groups that contribute most significantly to predictive performance.

%=======================================
% Feature Selection
%=======================================
\begin{table}[h!]
\centering
\small
\begin{tabular}{ll}
\toprule
Feature Set & Code  \\
\midrule
All features (full feature set) &  All   \\
No collinear features  &  NoCol\\
No Clinical Profile features  &  NoProf\\
No Comorbidities features & NoCom \\
No Neurological Examination features & NoNeuro \\
No Michigan Neuropathy Screening Instrument feature & NoMSI  \\
No Sudoscan features & NoSudo \\
Sudoscan features only & Sudo \\
Sudoscan and Clinical Profile features only & SudoProf \\
Sudoscan and Comorbidities features only & SudoCom \\
Sudoscan and Neurological Examination features only & SudoNeuro \\
Sudoscan and MNSI features only & SudoMnsi \\
\bottomrule
\end{tabular}
\caption{Feature sets examined during the baseline model evaluation phase. Twelve distinct feature sets were constructed to systematically assess the contribution of each clinical feature group — Clinical Profile, Comorbidities, Neurological Examination, MNSI, and Sudoscan — both individually and in combination. The "All" set serves as the full-feature reference, while ablation sets (e.g., NoCol, NoProf, NoCom) isolate the effect of removing one feature group at a time. Additive sets anchored on Sudoscan (e.g., SudoProf, SudoCom, SudoNeuro, SudoMnsi) were additionally constructed given Sudoscan's established role in peripheral neuropathy screening, allowing assessment of its complementary value when paired with other feature groups.}
\label{tab:feature_sets}
\end{table}

In addition, multicollinearity was addressed using the Variance Inflation Factor (VIF). Features exhibiting a VIF greater than 5 were considered highly collinear and were removed, resulting in a reduced feature set aimed at improving model stability and interpretability.

To further investigate the predictive value of specific diagnostic modalities, models were also trained exclusively on nerve conduction study (NCS) features and Sudoscan features, respectively. This allows for a focused evaluation of their standalone predictive capabilities and their potential role in clinical decision support.

\subsubsection{Hyperparameter Optimization}
Following the comparative evaluation, the best-performing model was selected based on the results of the previous steps. In this study, CatBoost demonstrated superior performance and was therefore chosen for further optimization, as detailed in the Results and Discussion section.

Hyperparameter optimization was conducted using a Repeated Stratified k-Fold framework with 3 splits. For each split, the training portion of the data was subjected to a further 3-fold Bayesian optimization, with ROC-AUC as the objective function. This approach efficiently explores the hyperparameter space by balancing exploration and exploitation, leading to improved model configurations.

The optimal hyperparameters obtained from each training fold were then used to generate predictions on the corresponding test fold. A classification threshold was subsequently selected for each split, allowing for the calibration of sensitivity and specificity in accordance with the study objectives.

The mean and standard deviations of the threshold, ROC-AUC, Youden index, specificity and sensitivity are reported for the 3 splits, as well as the confidence interval at 5\% for thr ROC-AUC and the Youden index.

This multi-stage optimization framework ensures that the final model is both well-tuned and rigorously validated, enhancing its reliability and generalizability.

\subsection{Counterfactual Generation}
%%%%%%%%%%%%%%%%%%%%%%%%%%%%%%%%%%%%%%%
To provide interpretable insights into the predictions of the trained machine learning models, counterfactual explanations were generated using the DiCE (Diverse Counterfactual Explanations) framework \cite{mothilal2020explaining}. DiCE is well-suited for this study due to its model-agnostic capability, enabling its application across a wide range of classifiers, including complex ensemble models. Furthermore, it supports diversity-aware counterfactual generation, allowing the identification of multiple plausible feature modifications that can alter model predictions—an important requirement in clinical decision support. Its ability to incorporate feature constraints (e.g., permitted ranges and immutable variables) ensures that generated counterfactuals remain realistic and clinically meaningful.

\subsubsection{Feature Suitability as Counterfactuals}
Table~\ref{tab:cf_features} shows the feature set and their suitability for counterfactual analysis. We note that nerve conduction studies and Sudoscan features are marked as \textit{conditional} counterfactuals. These are biomarker outcomes of neuropathy. While these are not directly actionable inputs, they can be treated as counterfactuals that answer "what must change". Intervention is left to the physician. 

% =======================================
% Table of Counterfactual Features
% =======================================
\begin{table}[h!]
\centering
\footnotesize
\begin{tabular}{p{0.1\linewidth} p{0.25\linewidth} p{0.1\linewidth} p{0.35\linewidth}}
\toprule
\textbf{Group} &  \textbf{Feature(s)} & \textbf{Suitable?} & \textbf{Reason} \\
\midrule
Profile      & Patient's sex (M or F) & No & Cannot be changed \\
             & Patient's age & No & Cannot be changed \\
             & Other symptoms & No & Too vague to have an intervention \\
             & Diabetes Mellitus duration & No & Cannot be reversed \\
             & Patient is on insulin  & Yes & Treatment can be started, stopped, or changed \\
             & Hemoglobin A1c & Yes & Modifiable through lifestyle, diet, and medication \\
\midrule
Commorbidity & Hypertension & Yes & Manageable through antihypertensives and lifestyle \\
             & Peripheral Arterial Occulsive Disease & Conditional & Partially treatable (revascularization, lifestyle) but not fully reversible \\
             & Dyslipidemia & Yes & Modifiable through statins, diet, exercise\\
             & Chronic Kidney Disease & Conditional & Partially modifiable (nephroprotective interventions) but not reversible \\
             & Guillain-Barre Syndrome & No & Acute/post-infectious, not a modifiable \\
\midrule
Neurology \newline Examination & All  & No &  These are measurements of an examination direct intervention for which is difficult to define. \\
\midrule
MNSI            & Michigan Neuropathy \newline  Screening Instrument & No & This is measurement of an examination direct intervention for which is difficult to define.   \\
\midrule
Sudoscan & All & No &  These are measurements of a procedure direct intervention for which is difficult to define.\\
\bottomrule
\end{tabular}
\caption{Actionability classification of input features for counterfactual explanation generation. Each feature was evaluated along three criteria: clinical modifiability, reversibility, and intervention feasibility. Features rated Suitable are both modifiable and clinically actionable (e.g., HbA1c via diet or medication; insulin use via prescription change). Features rated Conditional are partially manageable but not fully reversible (e.g., Chronic Kidney Disease through nephroprotective interventions; Peripheral Arterial Occlusive Disease through revascularization). Features rated Unsuitable are either immutable (sex, age, diabetes duration), too vague for targeted intervention (other symptoms), or reflective of examination or procedural measurements rather than modifiable patient states (neurological examination findings, MNSI score, Sudoscan readings). Only Suitable and Conditional features were included in counterfactual generation to ensure that all generated counterfactual scenarios correspond to clinically meaningful and potentially achievable patient states.}
\label{tab:cf_features}
\end{table}

\subsubsection{General Setup}
Here we discuss the global permitted range and configuration parameters to be used by the DiCE object.

\textit{Global Permitted Range}.
To ensure feasibility of generated counterfactuals, a global permitted range was defined for all continuous features. The bounds were computed as follows:
\begin{itemize}
    \itemsep0em 
    \item Minimum: $min(0,x_{min} - \sigma)$
    \item Maximum: $x_{max} + \sigma$
\end{itemize}
where $x_{min}$ and $x_{max}$ denote the observed minimum and maximum values, respectively, of each feature in the entire dataset, while $\sigma)$ is the standard deviation. This formulation prevents unrealistic values while allowing controlled extrapolation beyond the observed data.

\textit{DiCE Configuration Parameters}. The configuration of the DiCE framework is summarized in Table~\ref{tab:dice_setup}. The genetic algorithm was selected as the optimization method because of the presence of multiple categorical features which renders the use of gradients for optimization not feasible. Genetic algorithm is also effectiveness in exploring high-dimensional feature spaces and generating diverse counterfactuals under multiple constraints.

% =======================================
% DICE General Configuration
% =======================================
\begin{table}[h!]
\centering
\small
\begin{tabular}{lll}
\toprule
Component & Parameter & Value / Description \\
% \multirow{3}{*}{Data (dice\_ml.Data)} 
\midrule
Data (dice\_ml.Data) & continuous\_features & continuous features as defined in Table~\ref{tab:clinical_data}-A \\
& categorical\_features & binary categorical features as defined in Table~\ref{tab:clinical_data}-A \\
& permitted\_range & global permitted range as defined in the section \textit{General Setup}\\
\midrule
% \multirow{3}{*}{Model (dice\_ml.Model)} 
Model (dice\_ml.Model) & model & trained model under the mean threshold \\
& backend & sklearn \\
& model\_type & classifier \\
\midrule
DiCE object(dice\_ml.Dice)  & method & genetic \\
\bottomrule
\end{tabular}
\caption{Configuration parameters used to initialize the DiCE (Diverse Counterfactual Explanations) framework for post-hoc model interpretability. Three core objects were instantiated: the Data object (dice\_ml.Data), which defines the continuous and binary categorical input features alongside their globally permitted ranges; the Model object (dice\_ml.Model), which wraps the trained scikit-learn classifier operating under the mean decision threshold; and the DiCE object (dice\_ml.Dice), which executes counterfactual generation using a genetic algorithm. Together, these components enable the generation of actionable counterfactual explanations that identify the minimal feature changes required to alter a patient's predicted DPN classification.}
\label{tab:dice_setup}
\end{table}

\subsubsection{Local Counterfactual Analysis}
We describe here how instances were chosen for counterfactual generation and how the permitted range was computed for each instance.

\textit{Targeted Instance Selection}.
Local counterfactual explanations were generated from misclassified and borderline cases. Borderline cases are for instances whose predicted probabilities lie within ±0.2 of the mean decision threshold obtained during hyperparameter optimization. These instances represent cases where small feature changes may alter the predicted outcome, making them particularly informative for analysis.

\textit{Local Permitted Range}.
For each selected instance, an instance-specific permitted range was defined for each feature. Similar to the global permitted range, this ensures that generated counterfactuals remain both realistic and tailored to the specific instance.

\begin{itemize}
    \itemsep0em 
    \item Minimum: $min(0, instance\_value - \sigma)$
    \item Maximum: $ instance\_value + \sigma)$
\end{itemize}

\subsubsection{Local Counterfactual Generation Parameters}
The parameters used for generating local counterfactuals via $generate\_counterfactuals$ are summarized in Table~\ref{tab:local_cf}. The selected parameters balance proximity, diversity, and sparsity, ensuring that generated counterfactuals are close to the original instances, sufficiently varied, and involve minimal feature changes. This results in actionable and interpretable explanations suitable for clinical applications.

\begin{table}[h!]
\centering
\small
\begin{tabular}{ll}
\hline
\textbf{Parameter} & \textbf{Value / Description} \\
\hline
instances & Selected misclassified and borderline instances \\
total\_CFs & 20 \\
desired\_class & opposite \\
permitted\_range & local permitted range \\
features\_to\_vary & all features except those marked unsuitable in Table~\ref{tab:cf_features} \\
stopping\_threshold & mean threshold from hyperparameter optimization \\
posthoc\_sparsity\_algorithm & binary \\
posthoc\_sparsity\_param & 0.05 \\
proximity\_weight & 0.5 \\
diversity\_weight & 0.1 \\
categorical\_penalty & 0.1 \\
algorithm & DiverseCF \\
\hline
\end{tabular}
\caption{Parameters passed to the $generate\_counterfactuals$ function for local counterfactual generation. For each selected misclassified or borderline instance, up to 20 counterfactuals targeting the opposite class were generated, with feature variations constrained to locally permitted ranges.}
\label{tab:local_cf}
\end{table}

% \subsubsection{Global Feature Importance}
% Since running DiCE’s $global\_feature\_importance$ failed to complete within reasonable time, the global feature importance was computed manually per model as the ratio of total number of times a feature was changed in a model, to the total number of counterfactuals for that model. 

% Global feature importance was computed using DiCE’s $global\_feature\_importance$ method, which identifies features that are most frequently altered across counterfactuals, thereby indicating their influence on model predictions. The parameters used are summarized in Table~\ref{tab:global_cf}.

% \begin{table}[h!]
% \centering
% \small
% \begin{tabular}{ll}
% \toprule
% \textbf{Parameter} & \textbf{Value / Description} \\
% \midrule
% total\_CFs & 20 \\
% features\_to\_vary & INSULIN, HBA1C, HPN, PAOD, DSLPDMIA, CKD (see Table~\ref{tab:cf_features}) \\
% permitted\_range & global permitted range \\
% stopping\_threshold & mean threshold from hyperparameter optimization \\
% posthoc\_sparsity\_algorithm & binary \\
% \bottomrule
% \end{tabular}
% \caption{Global Counterfactual Importance Parameters. These are the parameters for computing the global importance with a call to \textit{global\_feature\_importance}.}
% \label{tab:global_cf}
% \end{table}


\section{Results and Discussion}
%%%%%%%%%%%%%%%%%%%%%%%%%%%%%%%%%

In this section, we discuss the performance of the machine learning models, the hyperparameter optimization and local explainability of features through feature importance and SHAP specifically for CatBoost, the selected ML model for counterfactual generation. We then proceed to an analysis of the counterfactuals by discussing the  global feature importances, describing local counterfactual instances per model split, and finally, providing metrics for the counterfactuals per split.

\subsection{Feature Set and Model Selection}

\subsubsection{Baseline Model Evaluation: Rationale and Metric Selection}
The primary objective of this baseline evaluation was to identify the most promising classification algorithm for subsequent optimization as a machine learning-based pre-diagnostic screening tool for Diabetic Peripheral Neuropathy (DPN). Fourteen algorithms spanning diverse learning paradigms were assessed using default hyperparameters and repeated stratified k-fold cross-validation (10 repeats × 4 folds; 40 total runs) to obtain stable, low-variance performance estimates without data leakage.
Given the inherent class imbalance in this dataset — 131 DPN-confirmed patients versus 59 unconfirmed cases (approximately 69\% positive prevalence) — AUPRC was chosen as the primary ranking metric in preference to ROC-AUC. This decision is well-supported in the methodological literature: ROC-AUC can present an overly optimistic picture of classifier performance under class imbalance because it incorporates the true negative rate, which inflates when the negative class is underrepresented. AUPRC, by contrast, focuses exclusively on the precision-recall trade-off and is therefore more sensitive to performance differences among classifiers on the minority class \cite{saito2015precision, davis2006relationship}. The Naïve baseline — which predicts the majority class uniformly — achieved an AUPRC of only 0.68 across all feature sets and conditions, despite attaining perfect sensitivity (1.0) and a reasonable accuracy of 68.4\%, illustrating precisely why accuracy and sensitivity alone are insufficient evaluation criteria in imbalanced screening tasks.

\subsubsection{Model Performance and Selection}
Across all 40 runs on the full feature set, ensemble tree-based methods consistently outperformed other model families (Table~\ref{tab:selection_metrics}). CatBoost and Extra Trees achieved the highest mean AUPRC (0.890), with a median AUPRC distribution that was both elevated and comparatively compact relative to other top models (Figure~\ref{fig:auprc_violins}), reflecting strong performance consistency across data partitions. Random Forest and RBF SVM also achieved a high AUPRC (0.889 and 0.887, respectively) but the widest distributional spread among top models — suggesting heightened sensitivity to the specific composition of training and test folds, which is a practical concern in small-sample clinical datasets.

From a clinical deployment standpoint, the balance between sensitivity and specificity is of critical importance in pre-diagnostic screening. A tool intended to triage patients toward the more costly and less accessible Nerve Conduction Study (NCS) must prioritize minimizing false negatives (i.e., maintaining high sensitivity) to avoid denying care to patients who truly have DPN. 

Under this clinical framing, CatBoost's sensitivity of 0.912 — combined with it achieving the highest AUPRC — is particularly compelling. RBF SVM achieved the highest sensitivity overall (0.939), but at the cost of the lower specificity in the cohort (0.5). CatBoost provides a more balanced operating point, and its gradient-boosting architecture with native categorical feature handling offers additional practical advantages in heterogeneous clinical tabular data \cite{dorogush2018catboost}.

\subsubsection{Feature Set Analysis}
The heatmap of average AUPRC across feature sets Figure~\ref{fig:auprc_heatmaps} reveals several findings with direct clinical relevance. First, the top models (CatBoost, Extra Trees, Random Forest, RBF SVM) are remarkably robust to the removal of most feature groups: ablating Clinical Profile, Comorbidities, Sudoscan or MNSI features produced negligible AUPRC degradation for these models ($\le 0.02$ points). This robustness implies that no single feature group is singularly indispensable, and that redundancy across the feature space is substantial. The exception are Neurological Examination features, the removal of which showed the greatest performance degradation suggesting that these features may be strong predictors. 

The Sudoscan only feature set understandably lowered the performance, while the additive Sudoscan-anchored feature sets (SudoProf, SudoCom, SudoNeuro, SudoMnsi) generally demonstrated modest but inconsistent improvements over the Sudoscan only set, suggesting diminishing returns from adding other feature groups once Sudoscan data are available. 
One exception worth noting is that the SudoNeuro combination (Sudoscan with Neurological Examination findings) maintained high performance across more model types, indicating that simple bedside neurological signs — vibration sensation, pinprick, light touch, and ankle reflex — may complement electrochemical skin conductance data in a synergistic manner that benefits less flexible classifiers.

We note that since these metrics reported are based on default parameters for the methods, a more complete picture of the importances of the features will be explored during the hyperparameter optimization that follows. 


\subsubsection{Reproducibility and Stability of Estimates}
We use of 10-repeat × 4-fold stratified cross-validation was chosen deliberately to obtain stable performance estimates given the modest sample size (n = 190). With a single k-fold split, estimates can exhibit high variance due to the stochastic partitioning of a small dataset \cite{kim2009estimating}. 

Repeated cross-validation substantially reduces this variance, and the violin plot distributions (Figure~\ref{fig:auprc_violins}) confirm that top models maintain consistent AUPRC distributions with limited lower-tail excursions, lending confidence to the relative ranking of algorithms. The absence of hyperparameter tuning or threshold optimization in this phase is also intentional: these results represent genuine baseline performance, ensuring that the subsequent model selection and optimization pipeline is not confounded by information leakage from the full dataset.

\subsubsection{Study Limitations at the Baseline Evaluation Stage}
Several limitations must be acknowledged at this stage. The dataset is relatively small (n = 190), and though repeated stratified cross-validation mitigates instability, external validation on an independent cohort remains necessary before clinical deployment. The class imbalance, while modest compared to many medical screening contexts, may still influence threshold-dependent metrics (sensitivity, specificity, accuracy) reported here, all of which used the default 0.5 probability threshold — a suboptimal choice for screening applications. Threshold calibration and optimization  will be addressed in the subsequent modeling phase. 

\subsubsection{Selected Feature Set and Model}
Because of CatBoost's notable performance (AUPRC of 0.89, narrow AUPRC distribution and sensitivity of 0.912), we choose it as the model to undergo the next stages of this research, namely, hyperparameter optimization and counterfactual generation. Moreover, since the full feature dataset as the data set that the ML model will be trained on since it consistently achieved high performance across multiple methods.

% selection Stage - All Metrics Table
% -----------------------------------
\begin{table}[h!]
\centering
\footnotesize
\begin{tabular}{lcccccccc}
\toprule
Algorithm & AUPRC & ROC-AUC & Sensitivity & Specificity & Accuracy & Precision & Youden Index & F1 Score\\
\midrule
CatBoost & 0.890 & 0.817 & 0.912 & 0.567 & 0.803 & 0.821 & 0.479 & 0.864 \\
Extra Trees & 0.890 & 0.807 & 0.879 & 0.548 & 0.775 & 0.810 & 0.428 & 0.842 \\
Random Forest & 0.889 & 0.817 & 0.902 & 0.568 & 0.797 & 0.820 & 0.471 & 0.858 \\
RBF SVM & 0.887 & 0.824 & 0.939 & 0.500 & 0.801 & 0.804 & 0.439 & 0.866 \\
Naive Bayes & 0.879 & 0.808 & 0.656 & 0.633 & 0.649 & 0.729 & 0.289 & 0.647 \\
LDA & 0.871 & 0.800 & 0.845 & 0.627 & 0.776 & 0.832 & 0.472 & 0.837 \\
Logistic Regression & 0.865 & 0.792 & 0.874 & 0.582 & 0.782 & 0.821 & 0.456 & 0.845 \\
XGBoost & 0.861 & 0.772 & 0.849 & 0.563 & 0.759 & 0.809 & 0.412 & 0.827 \\
Gradient Boosting & 0.858 & 0.766 & 0.863 & 0.558 & 0.767 & 0.811 & 0.422 & 0.834 \\
Linear SVM & 0.854 & 0.775 & 0.849 & 0.570 & 0.761 & 0.812 & 0.419 & 0.828 \\
LightGBM & 0.851 & 0.765 & 0.849 & 0.553 & 0.755 & 0.805 & 0.402 & 0.825 \\
kNN & 0.834 & 0.753 & 0.856 & 0.508 & 0.746 & 0.791 & 0.365 & 0.821 \\
SGDClassifier & 0.812 & 0.716 & 0.735 & 0.422 & 0.637 & 0.756 & 0.157 & 0.699 \\
Decision Tree & 0.758 & 0.648 & 0.775 & 0.520 & 0.695 & 0.778 & 0.295 & 0.775 \\
Naive & 0.684 & 0.500 & 1.000 & 0.000 & 0.684 & 0.684 & 0.000 & 0.812 \\
\bottomrule
\end{tabular}
\caption{Aggregate performance metrics — AUPRC, ROC-AUC, Sensitivity, Specificity, Accuracy, Precision, Youden Index, and F1 Score — for all 14 classification algorithms evaluated on the full feature set under 10-repeat × 4-fold stratified cross-validation (40 total runs; n = 190). All metrics are reported as averages across runs. CatBoost achieves the highest AUPRC (0.89) alongside the highest sensitivity (0.912), making it the most suitable candidate for a DPN pre-diagnostic screening tool where minimizing false negatives is clinically paramount. Extra Trees and Random Forest follow closely in AUPRC (both 0.89 and 0.889, respectively), with competitive sensitivity values. RBF SVM attains the highest overall sensitivity (0.939) at the cost of the lowest specificity (0.5), indicating a pronounced bias toward positive-class prediction. The Naïve classifier — predicting the majority class throughout — achieves perfect sensitivity and zero specificity, representing an uninformative lower bound for comparison. No threshold optimization or hyperparameter tuning was applied at this stage; all models used default parameters to provide an unbiased baseline for subsequent model selection.)
}
\label{tab:selection_metrics}
\end{table}

% AURPRC HEATMAPS
%------------------------------
\begin{figure*}[h]
    \centering
    \includegraphics[width=0.8\textwidth]{figures2/modeling/auprc_summary_table.png}
    \caption{Heatmap of average AUPRC scores across 14 classification models (rows) and 12 feature sets (columns), computed over 40 repeated stratified k-fold runs. Color intensity encodes AUPRC magnitude. The top models (CatBoost, Extra Trees, Random Forest, RBF SVM) are remarkably robust to the removal of most feature groups: ablating Clinical Profile, Comorbidities, Sudoscan or MNSI features produced negligible AUPRC degradation for these models. The NoNeuro column shows the greatest performance degradation across different models suggesting that these features may be strong predictors. We highlight that the SudoNeuro combination (Sudoscan with Neurological Examination findings) maintained high performance across more model types, achieving values close or sometimes better than the full feature set. The Naïve baseline remains uniformly at 0.68 across all feature sets, serving as a meaningful lower bound. Since these values are based on default parameters for the algorithms, a more complete picture of the importances of the features will be explored during the hyperparameter optimization that follows.}
    \label{fig:auprc_heatmaps}
\end{figure*}

% Violin Plots
%-------------
\begin{figure*}[h]
    \centering
    \includegraphics[width=0.7\textwidth]{figures2/modeling/auprc_violin.png}
    \label{fig:violin_auprc}
    \caption{Distribution of Area Under the Precision-Recall Curve (AUPRC) scores across 40 evaluation runs (10-repeat × 4-fold stratified cross-validation) for all 14 classification models evaluated on the full feature set. Each violin depicts the spread and density of AUPRC values, with the embedded box plot indicating the median (white line) and interquartile range. Models are ordered from left to right by descending mean AUPRC (see Table~\ref{tab:selection_metrics} for actual values). CatBoost and Extra Trees, the top 2 methods, demonstrate comparatively compact distributional spread reflecting both strong average performance and consistency across data partitions. Random Forest has a wider spread suggesting sensitivity to data partitioning.  In contrast, the Naïve baseline model exhibits minimal variance, consistent with its degenerate majority-class prediction strategy.}
    \label{fig:auprc_violins}
\end{figure*}

\subsection{CatBoost Model Training and Optimization}
In this subsection, we describe the model training and hyperparameter optimization CatBoost model for the full feature set (all data in Table~\ref{tab:clinical_data}-A). Four CatBoost models were independently trained and optimized across a 4-fold stratified cross-validation scheme. The test set of each fold underwent Bayesian optimization and threshold selection, under a 3-fold scheme (optimization of parameters for 2/3 of the fold's training set, threshold selection for 1/3 of the fold's training set). The optimization targeted the Area Under the Precision-Recall Curve (AUPRC) using the CatBoost parameters tree depth, learning rate, L2 leaf regularization, and positive class scale weight. The reported metrics were obtained through the hold out set of each outer fold, completely unseen by both the optimization and threshold selection process, thus preventing any data leakage. 

\subsubsection{Model Performance and Clinical Suitability for Pre-Diagnostic Screening}
% ============================================================================================

% CatBoost Performance Metrics First 4 Models
%--------------------------------------------
\begin{table}[t]
\centering
\footnotesize
\begin{tabular}{llcccccc}
\toprule
& & Model 0 & Model 1 & Model 2 & Model 3 & Mean & Std\\
\midrule
Parameter &  threshold & 0.427 & 0.392 & 0.447 & 0.470 & 0.434 & 0.026 \\
\midrule
Catboost Parameters & depth & 5.000 & 8.000 & 10.000 & 4.000 & 6.750 & 2.133 \\
& learning\_rate & 0.249 & 0.004 & 0.007 & 0.004 & 0.066 & 0.095 \\
& l2\_leaf\_reg & 2.713 & 9.273 & 7.789 & 9.909 & 7.421 & 2.527 \\
& scale\_pos\_weight & 0.885 & 0.906 & 0.748 & 0.988 & 0.882 & 0.077 \\
\midrule
Metrics & AUPRC & 0.925 & 0.857 & 0.892 & 0.907 & 0.895 & 0.022 \\
& ROC-AUC & 0.848 & 0.824 & 0.835 & 0.821 & 0.832 & 0.010 \\
& Sensitivity & 0.879 & 0.970 & 0.938 & 0.906 & 0.923 & 0.030 \\
& Specificity & 0.467 & 0.400 & 0.733 & 0.533 & 0.533 & 0.112 \\
& F1 Score & 0.829 & 0.865 & 0.909 & 0.853 & 0.864 & 0.026 \\
& Youden Index & 0.345 & 0.370 & 0.671 & 0.440 & 0.456 & 0.115 \\
& Precision & 0.784 & 0.780 & 0.882 & 0.806 & 0.813 & 0.037 \\
& Accuracy & 0.750 & 0.792 & 0.872 & 0.787 & 0.800 & 0.040 \\
\bottomrule
\end{tabular}
\caption{Best parameters and performance metrics of 4 CatBoost models independently trained and tested over 4 data splits. Optimized hyperparameters and classification performance metrics for four independently trained and evaluated CatBoost models across a 4-fold stratified cross-validation scheme. Each model was tuned using Bayesian optimization targeting the Area Under the Precision-Recall Curve (AUPRC). Threshold, tree depth, learning rate, L2 leaf regularization, and positive class scale weight were optimized per fold. Performance metrics are reported for the held-out test partition of each fold, together constituting evaluation across the entire dataset.}
\label{tab:catboost_metrics_first4}
\end{table}

% CATBOOST AUPRC
%------------------------------
\begin{figure*}[h]
\centering
\includegraphics[width=0.6\textwidth]{figures2/explainability/catboost_all_splits_auprc.png}
\caption{Precision-Recall (PR) curves for each of the four CatBoost models (Models 0–3) trained and evaluated on independent held-out test partitions across a 4-fold cross-validation scheme. Individual model Average Precision (AP) values were 0.93, 0.86, 0.89, and 0.90, respectively. The bold blue curve represents the mean PR performance across all folds (mean AP = 0.89 ± 0.03 SD). The red dashed horizontal line denotes the no-skill baseline (AP = 0.69), corresponding to the positive class prevalence in the study cohort (131/190 = 68.9\%). All four models substantially exceeded the no-skill baseline across the full recall range, demonstrating robust and consistent discriminative performance for identifying patients who should be referred for Nerve Conduction Study (NCS). AP was computed as the weighted mean of precisions across all classification thresholds.}
\label{fig:catboost_auprc}
\end{figure*}

% CatBoost Repeated k-fold Performance Metrics
%---------------------------------------------
\begin{table}[t]
\centering
\footnotesize
\begin{tabular}{llrrrr}
\toprule
&  & Mean & STD & CI Lower & CI Upper \\
\midrule
Parameter & Threshold & 0.363 & 0.136 & 0.320 & 0.405 \\
\midrule
Metrics &  AUPRC & 0.885 & 0.054 & 0.868 & 0.902 \\
& ROC-AUC & 0.810 & 0.064 & 0.790 & 0.830 \\
& Sensitivity & 0.924 & 0.063 & 0.904 & 0.944 \\
& Specificity & 0.475 & 0.144 & 0.430 & 0.520 \\
& Accuracy & 0.782 & 0.045 & 0.768 & 0.796 \\
& Precision & 0.795 & 0.040 & 0.782 & 0.808 \\
& F1-score & 0.853 & 0.032 & 0.843 & 0.863 \\
& F-beta at 1.25 & 0.867 & 0.035 & 0.856 & 0.878 \\
& F-beta at 1.5 & 0.878 & 0.039 & 0.866 & 0.890 \\
& F-beta at 1.75 & 0.887 & 0.043 & 0.873 & 0.900 \\
& F-beta at 2.0 & 0.893 & 0.046 & 0.879 & 0.908 \\
& Youden Index & 0.399 & 0.126 & 0.359 & 0.438 \\
\bottomrule
\end{tabular}


\caption{Aggregate performance metrics for the CatBoost model evaluated under repeated stratified k-fold cross-validation (10 repeats × 4 folds; 40 total evaluation runs; n = 190). Metrics are summarized as the mean, standard deviation (STD), and 95\% confidence interval bounds (CI Lower, CI Upper) across all runs. The F-beta family of metrics ($\beta$ = 1.25, 1.50, 1.75, and 2.0) supplements the standard F1-score by progressively weighting sensitivity over precision, reflecting the clinical priority of minimizing false negatives in a pre-diagnostic screening context. The Youden Index (sensitivity + specificity - 1) is reported as a threshold-independent summary of overall discriminative balance. This repeated cross-validation scheme was performed to obtain statistically stable, bias-reduced estimates of out-of-sample performance that are less sensitive to any particular random partition of the dataset. The optimized decision threshold is reported separately from performance metrics, as it represents the operating parameter selected per fold during Bayesian optimization rather than an outcome of model discrimination.}
\label{tab:rkfold}
\end{table}

Table \ref{tab:catboost_metrics_first4} shows the best parameters and performance metrics for the four models, together with their mean and standard deviation. Mean AUPRC and Sensitivity (0.895 and 0.923) are slightly higher than the baseline unoptimized CatBoost model (0.89 and 0.912, see Table~\ref{tab:selection_metrics}). The high mean sensitivity across folds indicates that the models correctly identified over 90\% of DPN-positive patients on average. For a pre-diagnostic screening instrument intended to triage patients toward the more resource-intensive and often geographically inaccessible Nerve Conduction Study (NCS), the clinical priority is high sensitivity — minimizing false negatives and ensuring DPN cases are not missed at the screening stage [3]. 
The mean sensitivity across folds is consistent with thresholds recommended for screening tools in conditions where undetected disease has significant downstream consequences, including progression to foot ulceration, amputation, or autonomic failure [4, 5].  

The comparatively lower mean specificity of 0.533 (range: 0.400–0.733) reflects the deliberate calibration of decision thresholds — which ranged from 0.392 to 0.470 — toward recall maximization, a known and clinically acceptable tradeoff in screening paradigms where the cost of a missed true positive outweighs the cost of an unnecessary referral [6].

Variability in Youden Index across folds (0.345 to 0.671) highlights that model 2 achieved a more balanced sensitivity-specificity operating point, potentially attributable to the particular composition of its training partition. Nevertheless, the consistency of high sensitivity across all four folds, each covering a non-overlapping test segment of the full cohort, supports the robustness and generalizability of the approach. The use of 4-fold stratified cross-validation — rather than a single train-test split — ensured that every patient contributed to model evaluation, reducing selection bias and providing a more reliable estimate of out-of-sample performance for a dataset of this size [7].

The F1-scores, ranging from 0.829 to 0.909 (mean: 0.864), and the AUPRC values across all four models remaining well above 0.85, are competitive with published machine learning models for DPN detection. Prior works employing logistic regression, support vector machines, and neural networks on comparable clinical datasets have reported AUPRC values in the 0.75–0.88 range [8, 9], while CatBoost and gradient boosting approaches have demonstrated particular advantages in handling mixed feature types and class imbalance without extensive preprocessing — properties especially beneficial in clinical tabular data [10, 11].

Figure \ref{fig:catboost_auprc} shows the Precision-Recall (PR) curve for each of the four CatBoost models trained and evaluated across independent folds. The plot shows that the models consistently achieved high discriminative performance, with mean AUPRC of 0.89 ± 0.03. Given the inherent class imbalance in the study cohort — where 68.9\% of patients (131/190) carried confirmed DPN — the AUPRC is the more informative summary metric, as it is sensitive to model behavior among the minority (DPN-negative) class and is not inflated by the abundance of true negatives [1, 2]. The mean AP of 0.89 represents a 29\% relative improvement over the no-skill baseline (AP = 0.69), indicating that the CatBoost substantially exceeded chance discrimination.

\subsubsection{Repeated Cross-Validation}
% =======================================================================================================
While the 4-fold stratified cross-validation results Table~\ref{tab:catboost_metrics_first4} established that each independently trained model achieved strong discriminative performance, the generalizability of those estimates depends on the assumption that the single partitioning was representative of the underlying data distribution — an assumption that becomes increasingly fragile in datasets of modest size [7]. To address this limitation, repeated stratified k-fold cross-validation was performed over 10 independent repetitions of the 4-fold scheme, yielding 40 total evaluation runs and producing the aggregate estimates presented in Table~\ref{tab:rkfold}.

The mean AUPRC across all 40 runs was 0.885 (95\% CI: 0.868–0.902, STD = 0.054), closely corroborating the single-repeat estimate of 0.89 ± 0.03 obtained from Table~\ref{tab:catboost_metrics_first4}, and confirming that the strong precision-recall performance observed previously was not an artifact of a favorable random partition. Critically, the tight confidence interval — spanning just 3.4 percentage points — indicates that model performance is highly reproducible across varying data splits, an important criterion for clinical tool reliability [7, 18]. 

The mean sensitivity of 0.924 (95\% CI: 0.904–0.944) — the most clinically critical metric for a pre-diagnostic screening instrument — was the best-constrained of all metrics, with a relatively narrow confidence interval spanning just 4 percentage points and a STD of 0.063. This stability indicates that the model's ability to correctly identify DPN-positive patients is robust to partition variability, a finding of direct clinical relevance: in a screening context where the consequence of a missed DPN case includes delayed intervention, progressive nerve damage, and increased risk of foot ulceration or lower-limb amputation, a consistently high and predictable sensitivity is more valuable than a high average sensitivity with wide variability [4, 5]. By contrast, specificity showed the greatest instability across runs (mean = 0.475, STD = 0.144, 95\% CI: 0.430–0.520), confirming that the model's ability to correctly exclude DPN-negative patients was more sensitive to partition composition. This reflects the relatively smaller size of the DPN-negative class (n = 59) and the deliberate calibration of decision thresholds toward recall; specificity variability of this magnitude is consistent with published benchmarks for imbalanced clinical screening models [1, 6].

The F-beta metric family provides a principled framework for evaluating the sensitivity-precision tradeoff at varying clinical cost assumptions [19]. The progressive increase in mean F-beta scores from F1 = 0.853 at $\beta$ = 1.0 to F2 = 0.893 at $\beta$ = 2.0 reflects the model's stronger performance when sensitivity is weighted more heavily than precision, which is the appropriate operating regime for a screening tool. Importantly, the confidence intervals for all F-beta metrics remained narrow (range of CI widths: 0.022–0.029), confirming stable performance regardless of the chosen $\beta$ weighting. Clinicians or health system administrators calibrating the model for deployment can therefore select a $\beta$ value reflecting their local cost tolerance — for example, choosing $\beta$ = 1.5 (mean = 0.878, 95\% CI: 0.866–0.890) if a moderate emphasis on sensitivity over precision is deemed appropriate — with confidence that the reported performance bounds will generalize reliably.

The mean Youden Index of 0.399 (95\% CI: 0.359–0.438, STD = 0.126) was the most variable metric in relative terms, consistent with its dependence on both sensitivity and specificity, the latter being inherently unstable in imbalanced settings. Nevertheless, a Youden Index above 0.35 across the full confidence interval confirms that the model operates meaningfully above chance discrimination at the selected thresholds [20]. The relatively wide CI of the Youden Index further underscores that composite balance metrics are more partition-sensitive than sensitivity-dominant metrics in this class distribution, reinforcing the clinical rationale for prioritizing AUPRC and sensitivity as primary evaluation criteria for this use case.

Taken together, the repeated cross-validation results establish that the CatBoost model's performance estimates are statistically stable and not the product of favorable data splitting. The concordance between single-repeat (Table~\ref{tab:catboost_metrics_first4} ) and repeated cross-validation (Table~\ref{tab:rkfold} ) estimates across all metrics — within one standard deviation in every case — supports the reliability of the single-fold models as representative of expected deployment performance. This is a prerequisite for advancing the model toward prospective clinical validation, and addresses a common methodological criticism of machine learning studies in clinical medicine that rely on single train-test splits or single cross-validation repetitions [18, 21].


\subsection{Feature Importance and Model Interpretability}
% =========================================================

% Feature Importances
% ---------------------
\begin{figure}
\centering
\includegraphics[width=0.85\textwidth]{figures2/explainability/catboost_all_splits_feature_importances.png}
\caption{
CatBoost-derived feature importance scores for the 22 input features across all four independently trained models (Models 0–3). Features are color-coded by clinical domain group: Sudoscan-derived electrochemical skin conductance parameters (orange), patient demographic and clinical profile variables (blue), comorbidity indicators (green), clinical neurological examination findings (red/pink), and the Michigan Neuropathy Screening Instrument (MNSI) subscores (purple). Importance scores reflect the mean absolute contribution of each feature to reducing prediction loss across all trees in the ensemble. Features are ranked by their average importance across folds.}
\label{fig:feature_importances}
\end{figure}

% SHAP
%------------------------------
\begin{figure}
\centering
\includegraphics[width=0.85\textwidth]{figures2/explainability/catboost_all_splits_shap.png}
\caption{
 Heatmap of mean absolute SHAP (SHapley Additive exPlanations) values for each of the 22 input features across the four independently trained CatBoost models (Models 0–3). Each cell represents the mean absolute SHAP value for that feature in the corresponding model, reflecting the average magnitude of that feature's contribution to individual predictions. Feature color coding follows the same clinical domain groupings as Figure \ref{fig:feature_importances}. Higher SHAP values indicate greater average influence on the model output regardless of direction. Features are ordered by their average mean absolute SHAP value across folds.
 }
\label{fig:shap}
\end{figure}

Across all four folds, CatBoost feature importances (Figure~\ref{fig:feature_importances}) and SHAP-based explanations (Figure~\ref{fig:shap}) converged on a consistent set of clinically meaningful predictors, reinforcing the model's interpretability and lending biological plausibility to its decision logic — a critical consideration for clinical acceptance of machine learning tools [12].

Two neurological examination findings — decreased ankle reflex (DEC\_AR) and decreased vibration sensation (DEC\_VS) — emerged as the most consistently dominant predictors across all models. DEC\_AR was the single most important feature overall, with importance scores ranging from 9.74 to 14.43 and mean absolute SHAP values of 0.072–0.076 across all four folds. DEC\_VS similarly ranked among the top five features in every fold (importance: 5.78–10.24; SHAP: 0.039–0.072). This is clinically expected and biologically plausible: loss of ankle deep tendon reflexes and impaired vibration perception are among the earliest and most reliable large fiber neurological signs of DPN, and both are incorporated as cardinal diagnostic criteria in established clinical and electrophysiological staging frameworks [4, 5, 17]. Their consistent dominance across all partitions confirms that the model is anchoring its discriminative reasoning on the most diagnostically informative clinical signs — a reassuring finding for clinical translation and acceptance.

Multiple Sudoscan-derived features collectively contributed substantial predictive weight, forming the second most influential feature group in aggregate. Feet asymmetry percentage (FEET\_PCT\_ASYM) ranked consistently among the top features across all folds (importance: 7.66–9.88; SHAP: 0.018–0.048), as did the neuropathy score (NS; importance: 7.43–9.65; SHAP: 0.021–0.047) and the cardiac autonomic score (CAS; importance: 3.52–8.78; SHAP: 0.010–0.049). Foot-level mean electrochemical skin conductance (FEET\_MEAN\_ESC) and its hand-level counterpart (HAND\_MEAN\_ESC, HAND\_PCT\_ASYM) provided additional moderate contributions. Sudoscan measures small fiber sudomotor function through galvanic skin response and has been validated as a rapid, noninvasive surrogate marker of sudomotor neuropathy — an early manifestation of diabetic small fiber neuropathy that frequently precedes or accompanies the large fiber involvement detectable by NCS [13, 14]. The co-prominence of both large fiber neurological examination findings and multiple Sudoscan-derived small fiber parameters across all folds is particularly noteworthy: it suggests the model is integrating complementary information spanning the full spectrum of peripheral nerve fiber subtypes affected in DPN, rather than relying exclusively on any single modality. This multimodal integration may contribute to the model's robust sensitivity, as patients presenting with predominantly small fiber or predominantly large fiber involvement would each be captured through different feature pathways. This may also explain the initial finding on why, in the initial study of unoptimized models, the combination of the two set of features produced performance metrics close to the full set features (Figure~\ref{fig:auprc_heatmaps}).

Diabetes duration (DM\_DUR) ranked second in importance in Model 0 and remained consistently in the top five features across all folds (importance: 6.97–12.27; SHAP: 0.016–0.079), corroborating its well-established epidemiological role as a major risk factor for DPN. The relationship between disease duration and neuropathy risk is thought to reflect cumulative metabolic insult, oxidative stress, and neuroinflammation over time [4, 5]. Similarly, age (AGE), HbA1c (HBA1C), and insulin use (INSULIN) contributed consistently, albeit at lower magnitudes, consistent with their documented roles in DPN pathophysiology [4].

The MNSI score, reflecting structured patient-reported and clinician-assessed symptom burden, contributed meaningfully — particularly in Models 0 and 3 — with SHAP values of 0.010–0.041, suggesting it retains discriminative value even in a multivariate, data-driven framework alongside the more objective neurological examination and Sudoscan findings. This is relevant for settings where Sudoscan devices may not be available, as it implies the model draws meaningful signal from symptom-based instruments as well [15].

In contrast, comorbidities Hypertension (HPN) and Chronic Kidney Disease (CKD) received little feature importance and SHAP values, while peripheral arterial occlusive disease (PAOD) and Guillain-Barré syndrome (GBS) received near-zero importance in all four models, suggesting these comorbidities do not independently contribute to DPN prediction in this cohort — a finding that likely reflects their low prevalence in the study population and their pathophysiological distinction from diabetic neuropathy.

\subsection{Counterfactual Analysis}
%===============================================================
To move beyond predictive accuracy and toward clinical utility, this study employed counterfactual explanation analysis using the DiCE (Diverse Counterfactual Explanations) framework \cite{mothilal2020}. Counterfactual explanations address the question: "What is the minimal change to a patient's features that would result in a different model prediction?" Unlike global feature importance methods such as SHAP or LIME, which describe average model behavior, counterfactuals are inherently patient-specific and action-oriented, making them particularly suited to clinical decision support contexts \cite{wachter2017}.

% Local CF Summary Across Models
% -------------------------------
\begin{table}[h]
\footnotesize
\centering
\begin{tabular}{ccccccccc}
\toprule
Model & Candidates & Misclassified & Borderline & Instances with CF & Total CFs & Sparsity Mean & L1 Mean & L2 Mean \\
\midrule
0 & 12 & 12 & 2 & 4 & 58 & 4.19 & 4.17 & 3.29 \\
1 & 16 & 9 & 10 & 3 & 22 & 3.40 & 1.85 & 1.33 \\
2 & 11 & 7 & 7 & 2 & 14 & 3.32 & 1.18 & 0.82 \\
3 & 14 & 9 & 8 & 1 & 3 & 4.00 & 1.77 & 1.14 \\
\bottomrule
\end{tabular}
\caption{Aggregate summary of local counterfactual generation across four CatBoost models trained on cross-validation folds. Candidates denotes the number of test-set patients selected for counterfactual analysis based on being either misclassified or borderline (prediction probability within a defined low-confidence margin). Misclassified indicates the subset with incorrect predictions, while Borderline indicates low-confidence predictions (whether predictions were correct or not). Instances with CF reports how many candidate patients successfully yielded at least one valid counterfactual. Total CFs indicate the total number of counterfactuals the model produced for all instances. Sparsity Mean reflects the average number of features altered per counterfactual; L1 and L2 means quantify the magnitude of feature changes using Manhattan and Euclidean distances, respectively. Model 2 produced the most parsimonious counterfactuals (sparsity: 3.32; L1: 1.18; L2: 0.82), suggesting that its decision boundary permits smaller, more targeted feature perturbations to achieve a prediction reversal. Model 0 required the largest perturbations overall (L1: 4.17; L2: 3.29), indicating a more complex or less linearly separable boundary for its fold. The relatively low sparsity values across all models (range: 3.32–4.19 out of 6 actionable features) suggest that prediction reversals generally require changes to more than half of the available actionable features, highlighting the multifactorial nature of DPN risk.
}
\label{tab:localcf-model-level}
\end{table}

% Local CF Details Across Models
% -------------------------------
\begin{table}[h]
\footnotesize
\centering
\begin{tabular}{ccccccccc}
\toprule
Model & Patient ID & Probability & Margin & Outcome & CF Count & Sparsity Mean & L1 mean & L2 Mean \\
\midrule
0 & 60 & 0.52 & 0.10 & False Positive  & 32 & 3.33 & 1.23 & 1.15 \\
0 & 85 & 0.73 & 0.30 & False Positive  & 3 & 4.25 & 2.00 & 1.17 \\
0 & 99 & 0.22 & 0.20 & False Negative & 1 & 4.50 & 2.35 & 1.06 \\
0 & 135 & 0.80 & 0.37 & False Positive  & 22 & 4.70 & 11.09 & 9.78 \\
1 & 30 & 0.44 & 0.04 & False Positive  & 20 & 3.19 & 1.90 & 1.77 \\
1 & 58 & 0.46 & 0.07 & True Positive & 1 & 3.50 & 2.20 & 1.39 \\
1 & 74 & 0.46 & 0.06 & False Positive  & 1 & 3.50 & 1.45 & 0.84 \\
2 & 48 & 0.47 & 0.02 & True Positive & 13 & 3.14 & 0.86 & 0.78 \\
2 & 55 & 0.50 & 0.06 & True Positive & 1 & 3.50 & 1.50 & 0.87 \\
3 & 97 & 0.52 & 0.05 & True Positive & 3 & 4.00 & 1.77 & 1.14 \\
\bottomrule
\end{tabular}
\caption{Instance-level counterfactual statistics for individual candidate patients across four CatBoost models. Each row represents a patient for which counterfactuals were successfully generated. The model indicates where this patient was included as part of the test set.   
The rest of the columns report the model's predicted probability for the patient, the margin (distance from the 0.5 decision threshold), prediction outcome, the number of valid counterfactuals produced (CF Count), and the mean sparsity (number of features changed), L1 (Manhattan), and L2 (Euclidean) distances of those counterfactuals. Patients with very low margins (e.g., ID 48, margin = 0.02; ID 97, margin = 0.05) sit near the decision boundary, requiring only small feature changes (L1 $\le$ 1.77) to reverse classification. In contrast, Patient 135, a False Positive with probability 0.80 and margin 0.37, required substantially larger perturbations (L1 mean: 11.09; L2 mean: 9.78) across 22 counterfactuals, indicating a deeply embedded misclassification. Counterfactuals were generated for both misclassified (False Positive, False Negative) and correctly classified borderline (True Positive) cases, providing explanatory value for model uncertainty regardless of prediction correctness. The inclusion of True Positive borderline cases is clinically relevant: these represent patients correctly identified as high-risk but with low model confidence, for whom counterfactual scenarios can guide preventive clinical decisions.
}
\label{tab:localcf-listing}
\end{table}


\paragraph{Feature Actionability and the Counterfactual Candidate Pool}
A prerequisite for clinically meaningful counterfactuals is that the features being altered must correspond to interventions that are realistically achievable. Accordingly, features were stratified by actionability prior to counterfactual generation (Table~\ref{tab:cf_features}). Immutable demographic attributes (sex, age, DM duration) and examination-derived measurements (neurological findings, MNSI, Sudoscan readings) were excluded, as these do not correspond to modifiable patient states. This exclusion is not merely a technical constraint: counterfactuals generated from non-actionable features risk producing explanations that are statistically valid but clinically misleading [CITE:karimi2020]. The final actionable feature set—HbA1c, insulin use, hypertension, dyslipidemia, CKD, and PAOD—reflects a pragmatic and evidence-aligned subset of DPN risk factors for which clinical pathways exist. The inclusion of CKD and PAOD as conditional features acknowledges that while these conditions are not reversible, their management through nephroprotective and revascularization strategies can meaningfully alter the patient's risk trajectory.

\paragraph{Candidate Instances and Counterfactual Generation}
We identified misclassified instances (incorrect predictions) and borderline instances (prediction probability within 0.1) as candidate instances for counterfactual generation. A total of 53 candidates were identified but only 10 yielded counterfactuals across the 4 models after a predefined generation timeout of 1 hour. Table~\ref{tab:localcf-model-level} provides a summary of the counterfactuals generated per model while Table~\ref{tab:localcf-listing} provides details of the 10 candidates that successfully generated counterfactuals.

\subsubsection{Instance Level Counterfactual Analysis}
% =====================================================
In this section we provide an analysis and synthesis of the instances that were able to produce counterfactuals. Due to limitation of space, we will confine ourselves to the analysis of counterfactuals from Model 0. A feature change heatmap for these instances can be found in Figures~\ref{fig:localcf-m0-60}-\ref{fig:localcf-m0-135}. 

We chose Model 0 for detailed discussion since it produced the most counterfactuals and offers diverse scenarios covering borderline correct classification, confident misclassifications, and the sole example of a false negative outcome.  Collectively, these cases provide a rich and internally coherent account of Model 0's learned decision structure, its characteristic counterfactual behavior, and its specific strengths and vulnerabilities as a component of the ensemble pre-diagnostic screening framework. The counterfactuals for the Models 1-3 will be discussed in the aggregation-level counterfactual analysis that follows.

Model 0 applies a classification threshold of 0.4266 — a relatively moderate threshold that is neither the most liberal (Model 1 at 0.3917) nor the most conservative (Model 3 at 0.47). Four instances generated a total of 58 counterfactuals under this model. We proceed with the discussion of these instances and their counterfactuals below.


\begin{figure}[p]
\centering
\includegraphics[width=0.97\textwidth]{figures2/counterfactuals/060catboost_split0_local_cf.png}
\caption{Counterfactuals for Patient 60}
\label{fig:localcf-m0-60}
\end{figure}

\begin{figure}[p]
\centering
\includegraphics[width=0.97\textwidth]{figures2/counterfactuals/catboost_split0_local_cf_85.png}
\caption{Counterfactuals for Patient 85}
\label{fig:localcf-m0-85}
\end{figure}

\begin{figure}[p]
\centering
\includegraphics[width=0.97\textwidth]{figures2/counterfactuals/catboost_split0_local_cf_99.png}
\caption{Counterfactuals for Patient 99}
\label{fig:localcf-m0-99}
\end{figure}

\begin{figure}[p]
\centering
\includegraphics[width=0.97\textwidth]{figures2/counterfactuals/135catboost_split0_local_cf.png}
\caption{Counterfactuals for Patient 135}
\label{fig:localcf-m0-135}
\end{figure}


\paragraph{Patient 99 --- Confident False Negative}
\textit{Patient Profile}: Female, 42 years old; HbA1c 8.6\%; no insulin use; hypertension present; dyslipidemia present; no PAOD or CKD. Incorrectly predicted as Unconfirmed DPN at a probability of 0.2249, substantially below the Model 0 threshold of 0.4266 (margin = -0.2016), despite a confirmed clinical diagnosis.

Patient 99 represents the most clinically consequential misclassification type in a pre-diagnostic screening framework — a false negative — in which a patient with confirmed DPN was classified by Model 0 as Unconfirmed DPN, with a predicted probability of only 0.2249. The large negative decision margin of 0.2016  indicates that this patient was positioned far from Model 0's decision boundary of 0.4266, deep within the unconfirmed DPN classification region of the model's learned feature space. This deep mispositioning is particularly notable given that Patient 99 is the only false negative instance in Model 0.   

In the context of the proposed screening system, a false negative of this magnitude carries direct clinical risk: the patient would not be flagged for nerve conduction study referral, potentially delaying diagnosis and appropriate neuropathic management in a relatively young individual. The counterfactual heatmap (Figure~\ref{fig:localcf-m0-99}) underscores the severity of this misclassification: only a single counterfactual instance was generated.  

The analysis of this counterfactual must take into consideration that the suggestions are to flip the false negative prediction from Unconfirmed DPN to Confirmed DPN. The sole counterfactual proposed the simultaneous initiation of insulin therapy (INSULIN: 0→1), an increase in HbA1c (+0.7), the management of hypertension and dyslipidemia (HPN: 1→0; DSLPDMIA: 1→0), and the introduction of CKD (CKD: 0→1). While three of these five proposed changes are non-actionable — the removal of existing comorbidities and the introduction of a new pathological condition — and one (the HbA1c increase) is counterintuitive from a clinical management perspective, they are intuitive when understood from the perspective of flipping this patient's classification from Unconfirmed DPN to Confirmed DPN, that is, it can be understood that an increase in glycemic index, the introduction of chronic kidney disease and the recommendation to start taking insulin are characteristics, according to this model, of someone with Confirmed DPN. 

From a pre-diagnostic screening perspective, the false negative classification of Patient 99 serves as a critical reminder that even within a well-performing cross-validated model, individual fold-specific representational gaps can produce clinically consequential errors for atypical patient phenotypes, and that clinician judgment must remain an indispensable adjunct to model-based decision support in such cases.

\paragraph{Patient 60 --- Borderline False Positive}
\textit{Patient Profile}: Male, 55 years old; HbA1c 9.3\%; no insulin use; hypertension present; no dyslipidemia, PAOD, or CKD . Predicted as Confirmed DPN at a probability of 0.5216, marginally exceeding the Model 0 classification threshold of 0.4266 (margin = 0.095).

Patient 60 is a 55-year-old male classified by Model 0 as Confirmed DPN (probability = 0.5216) despite an unconfirmed clinical diagnosis.  The counterfactual heatmap (Figure~\ref{fig:localcf-m0-60}) revealed that glycated hemoglobin (HbA1c) was the single most decisive modifiable feature driving this patient's classification. Across 32 generated counterfactual instances (the largest counterfactual set produced by any patient in this analysis)  the predominant pathway to reclassification as Unconfirmed DPN consistently required a reduction in HbA1c, with proposed decrements ranging from -0.1 to -2.3 units. 

% The volume and granularity of this counterfactual set are characteristic of Model 0's behavior in the vicinity of its decision boundary: the model trained on this fold has learned a decision surface that is highly sensitive to incremental HbA1c variation for patients with this specific profile. 

Beyond glycemic control, a subset of counterfactuals (CF 9, CF 11–12, CF 15–16, CF 19–23, CF 28, CF 30, CF 32) proposed the introduction of comorbidities --- specifically dyslipidemia, CKD, or PAOD — as alternative or adjunctive reclassification pathways. The appearance of CKD as a proposed introduction in several of these counterfactuals is consistent with a pattern observed in other Model 0 cases, particularly Patient 85, suggesting that Model 0's training partition encoded a specific co-occurrence relationship between CKD and confirmed DPN classification that influences the geometry of its decision boundary. These counterfactuals — which suggest that a patient develop new diseases such as CKD or PAOD as a way to change the model's prediction — are not clinically useful, since no real intervention can or should induce a disease state in a patient. They are therefore labeled as non-actionable. One exception worth noting is CF 14, which proposed starting insulin therapy alongside a small reduction in HbA1c of −0.3. While this combination is closer to a real clinical scenario — since insulin initiation is a legitimate therapeutic decision — it cannot be considered a straightforward recommendation on its own, as the decision to start insulin depends on a broader clinical assessment beyond neuropathy risk alone. Overall, the counterfactual analysis for Patient 60 shows that the model's classification was not random or clinically unreasonable. The patient's elevated HbA1c was the main factor pushing the prediction toward confirmed DPN, and bringing it down through better diabetes management — such as medication adjustment, lifestyle modification, or closer glycemic monitoring — remains the most practical and meaningful way to potentially shift this patient's risk classification under Model Split 0.

\paragraph{Patient 135 --- Confident False Positive}
\textit{Patient Profile}: Female, 64 years old; HbA1c 6.7\%; no insulin use; hypertension present; dyslipidemia present; no PAOD or CKD. Predicted as Confirmed DPN at a notably high probability of 0.7971, substantially above the Model 0 threshold of 0.4266 (margin = 0.3705), despite an unconfirmed clinical diagnosis.

Patient 135 represents a diagnostically important false positive case, distinguished from the borderline case of Patient 60 (also a Model 0 case) by its considerably higher predicted probability (0.7971) and wider decision margin (0.3705). This within-split contrast is instructive: both patients were evaluated under the same decision structure by the model, yet they produced markedly different counterfactual landscapes — Patient 60 generating 32 graduated, predominantly HbA1c-driven counterfactuals and Patient 135 generating 22 counterfactuals of greater complexity and required perturbation magnitude (sparsity mean:4.70; L1 mean: 11.09; L2 mean: 9.78 --- see Table~\ref{tab:localcf-listing}). This divergence reflects the significant difference in their decision margins: Patient 60's boundary-adjacent positioning allowed Model 0 to generate numerous granular reclassification pathways, while Patient 135's deeper positioning within the confirmed DPN region — despite sharing the same threshold — required substantially larger and more complex feature perturbations to traverse the decision surface. The model's high-confidence prediction of 0.7971 for Patient 135 suggests that this patient's feature profile aligns strongly with the confirmed DPN prototype learned by Model 0's training partition, a characterization that warrants close clinical scrutiny given the unconfirmed diagnosis. Notably, the patient's HbA1c of 6.7\% falls within a range considered near-target for many diabetes management guidelines, rendering this case atypical with respect to the model's high-confidence prediction.

The most feasible pathways incorporating a potentially real therapeutic intervention is represented by those with moderate reduction in glycemic index as represented by, first, CF 11, CF 21, and CF 22, which require only moderate HbA1c reductions (-1.5 to -1.6) without altering the other features; second, CF 17, which require moderate HbA1c reduction of -1 and introducing insulin; and finally, CF 8,10 and 16 - which require moderate HbA1c reduction of -0.9 to -1 while addressing dyslipidemia and hypertension. Other counterfactuals suggested the introduction of comorbidities PAOD or CKD (CF 1-5, CF 7, CF 9, CF 14-17, CF 20). As previously explained, these suggestions are deemed unactionable. 

Collectively, these findings indicate that Model 0 has positioned Patient 135 deep within its confirmed DPN classification region. From a pre-diagnostic screening standpoint, this high-confidence false positive may represent a clinically meaningful risk signal warranting longitudinal monitoring or early nerve conduction study referral, particularly given the pronounced model confidence despite apparently near-controlled glycemia.

\paragraph{Patient 85 --- Confident False Positive}
\textit{Patient Profile}: Female, 69 years old; HbA1c 8.8\%; no insulin use; hypertension present; dyslipidemia present; no PAOD or CKD. Predicted as Confirmed DPN at a probability of 0.7269, substantially exceeding the classification threshold of 0.4266 (margin = 0.3003), despite an unconfirmed clinical diagnosis. 

Patient 85 represents a high-confidence false positive, with a predicted probability of 0.7269 placing this patient deep within the confirmed DPN classification region of the model's decision space. 
% The wide decision margin of 0.3003 — combined with the patient's advanced age (69 years), moderately elevated HbA1c (8.8\%), hypertension, and dyslipidemia — indicates that the model has encoded a cumulative cardiometabolic and age-related neuropathic risk profile strongly consistent with its learned representation of confirmed DPN cases. 
As with Patient 135, who presented a similarly high-confidence false positive under the same Model 0, the unconfirmed clinical diagnosis does not straightforwardly indicate a model error, but may instead reflect subclinical neuropathic involvement not yet captured by the diagnostic criteria applied.

The counterfactual heatmap (Figure~\ref{fig:localcf-m0-85}) for Patient 85 is notably sparse, generating only three counterfactual instances — a structural characteristic consistent with the model's deep positioning of this patient within the confirmed DPN region. Critically, all three counterfactuals were entirely non-actionable, suggesting the introduction of CKD (CKD: 0→1).

It is particularly striking that CKD introduction appeared in all three counterfactuals generated for this patient. In plain terms, the model could not find any realistic way to shift its prediction away from confirmed DPN without first imagining that this patient also had CKD — a condition he does not actually have. This pattern likely has one of two explanations. The first is that the model picked up a coincidental association in the training data, where patients without CKD happened to more frequently receive a confirmed DPN classification, causing the model to treat CKD absence as a subtle risk signal even though this relationship may not reflect a true biological mechanism. The second, more clinically interesting explanation, is that this association is real — meaning that within the patient population used to train the model, the specific combination of metabolic risk factors present in this patient genuinely tended to follow a different neuropathic risk course when CKD was also present, and the model learned to reflect this distinction in its predictions. Either way, because CKD introduction cannot be proposed as a clinical intervention, these counterfactuals remain non-actionable regardless of the underlying reason for their appearance.

Other suggestions from the counterfactuals follow a more intuitive logic for shifting this patient's classification toward Unconfirmed DPN. CF 1 and CF 3 both imagined a scenario where the patient no longer had hypertension, combined with a modest reduction in HbA1c of −0.6 and −0.7 respectively — in other words, the model suggests that if this patient's blood pressure condition were resolved and his blood sugar brought slightly closer to target, the prediction would change. CF 2 took a different approach, proposing that starting insulin therapy alongside a similar HbA1c reduction of −0.7 could achieve the same shift — suggesting that transitioning to insulin as part of an intensified diabetes management plan, together with modest glycemic improvement, represents an alternative pathway that the model recognizes as meaningfully lowering the patient's neuropathic risk profile. 

The low number of counterfactuals produced for Patient 85, together with the suggestions of modifications for other features show that high-confidence predictions in deeply positioned cases resist reclassification through any clinically feasible single-feature modification, underscoring the appropriateness of prioritizing such patients for urgent nerve conduction study referral within the proposed screening framework.

\paragraph{Synthesis of Local Counterfactual Instance Analyses}

Despite sharing a common classification threshold of 0.4266, the four Model 0 patients exhibited a remarkably wide range of predicted probabilities and decision margins, spanning from a large negative margin of −0.2016 (Patient 99, probability = 0.2249) to a large positive margin of 0.3705 (Patient 135, probability = 0.7971), with Patients 60 and 85 occupying intermediate positions of 0.095 (probability = 0.5216) and 0.3003 (probability = 0.7269) respectively. This wide intra-split dispersion in decision margins — covering nearly the full range of the model's predicted probability space — indicates that Model 0's decision surface is highly discriminating across a broad spectrum of metabolic and comorbidity profiles, positioning patients at vastly different distances from its threshold depending on the specific combination of features presented. This discriminative range is a favorable property of the model from a screening perspective, as it implies that the Model 0 threshold does not produce a dense clustering of cases near the boundary — which would indicate poor calibration — but rather a meaningful spread of confidence levels across the patient cohort.

The clinical implications of this margin heterogeneity are direct and actionable. Patient 135, positioned furthest above the threshold with a margin of 0.3705, represents the highest-priority candidate for nerve conduction study referral among the four Model 0 cases, as the model's high-confidence prediction reflects a robust alignment between this patient's feature profile and Model 0's learned confirmed DPN prototype. Patient 85, with a margin of 0.3003, occupies a similarly high-priority referral position. Patient 60, with a narrow positive margin of 0.095, is appropriately considered a moderate-priority case warranting watchful waiting and glycemic reassessment, given that modest improvements in HbA1c could legitimately shift the classification. Patient 99, with a large negative margin of −0.2016, represents the inverse scenario: a patient whose feature profile the Model 0 positions far from its confirmed DPN region, yet who carries a confirmed clinical diagnosis — the most clinically consequential and methodologically important case in the Model 0 cohort.

The collective behavior of Model 0 across its four analyzed cases carries several direct implications for the clinical deployment and continued development of this pre-diagnostic screening system. First, the Model 0 threshold of 0.4266 appears appropriately calibrated for high-sensitivity detection of confirmed DPN in patients with elevated comorbidity burden and moderate-to-high HbA1c — the dominant phenotype in its training partition — but may require threshold adjustment or ensemble weighting modification to improve sensitivity for younger, non-insulin-dependent confirmed DPN patients with moderate glycemic profiles. Second, the systematic appearance of CKD as a counterfactual pathway component across three of four Model 0 cases argues for an examination of CKD distribution and its relationship to DPN confirmation status within Model 0's training partition, with targeted data augmentation if a systematic imbalance is identified. Third, the decision margin values generated by Model 0 — spanning from a confident false negative to a highly confident false positive — suggest that this fold's outputs can meaningfully inform referral prioritization, with high-margin positive predictions (Patients 85 and 135) warranting urgent referral and narrow-margin positive predictions (Patient 60) supporting a monitored glycemic management pathway. Finally, the false negative case of Patient 99 serves as a specific and actionable quality indicator for Model 0, identifying a confirmed DPN phenotype for which this fold's representational capacity is demonstrably limited and for which ensemble-level classification — integrating predictions across all four splits rather than relying on any single fold — is particularly essential to maintain adequate screening sensitivity.

\subsubsection{Aggregate-level Counterfactual Analysis}
% =====================================================
At this point we analyze counterfactuals for all models on the aggregate level. This allows us to see across all models what features are being targeted to change the models' decisions, allowing us to assess the importance of these features and their possible interactions as well as determine possible clinical implications. Table~\ref{tab:localcf-changed} shows the feature-level frequency of changes across all counterfactuals (97) generated by 10 instances from 4 models.

% Global Importance Across Models
% -------------------------------

%  Local CF Change Features Across Models
% ---------------------------------------
\begin{table}[h]
\footnotesize
\centering
\begin{tabular}{ccccccccc}
\toprule
Model & Patient ID & CF Count & INSULIN & HBA1C & HPN & PAOD & DSLPDMIA & CKD \\
\midrule
0 & 60 & 32 & 1 & 29 & 0 & 2 & 5 & 7 \\
0 & 85 & 3 & 1 & 3 & 2 & 0 & 0 & 3 \\
0 & 99 & 1 & 1 & 1 & 1 & 0 & 1 & 1 \\
0 & 135 & 22 & 2 & 22 & 4 & 6 & 2 & 9 \\
1 & 30 & 20 & 0 & 20 & 0 & 0 & 5 & 0 \\
1 & 58 & 1 & 1 & 1 & 0 & 0 & 1 & 0 \\
1 & 74 & 1 & 1 & 1 & 0 & 0 & 1 & 0 \\
2 & 48 & 13 & 0 & 13 & 0 & 0 & 0 & 3 \\
2 & 55 & 1 & 1 & 0 & 1 & 0 & 1 & 0 \\
3 & 97 & 3 & 3 & 1 & 1 & 0 & 3 & 0 \\
\midrule
 & Totals & 97 & 11 & 91 & 9 & 8 & 19 & 23 \\
\bottomrule
\end{tabular}
\caption{Feature-level frequency of changes across all 97 counterfactuals generated from four CatBoost models. Each cell reports the number of counterfactuals for a given patient–model pair in which the corresponding feature was altered. Column totals summarize aggregate feature modification frequency across all instances. HbA1c (HBA1C) was modified in 91 of 97 counterfactuals (93.8\%), establishing it as the overwhelmingly dominant counterfactual driver. Chronic Kidney Disease (CKD) was the second most frequently modified feature (23 instances, 23.7\%), followed by Dyslipidemia (DSLPDMIA; 19 instances, 19.6\%), and Insulin use (INSULIN; 11 instances, 11.3\%). Hypertension (HPN) and Peripheral Arterial Occlusive Disease (PAOD) contributed minimally (9 and 8 instances, respectively). The concentration of changes in HbA1c across nearly all instances—regardless of model or patient—underscores that glycemic optimization represents the single most impactful actionable intervention for modifying DPN risk classification in this cohort. 
}
\label{tab:localcf-changed}
\end{table}

% Don't include as this could lead to critique; 
% Analysis from the local cf table is more meaningful
% \begin{figure}
% \centering
% \includegraphics[width=0.6\textwidth]{figures2/counterfactuals/global_cf.png}
% \caption{Global counterfactual feature importance across four independently trained CatBoost models (Models 0–3), derived from DiCE-generated counterfactuals on the hold-out test set of each cross-validation fold. Each cell represents the normalized proportion of counterfactuals in which that feature was altered to reverse the model's prediction, scaled from 0.0 (never changed) to 1.0 (changed in all counterfactuals). HbA1c (HBA1C) was the dominant feature across Models 0, 1, and 2 (importance scores: 0.95, 1.00, and 0.93, respectively), reflecting its central role in DPN risk modification. Model 3 exhibited a divergent pattern, with DSLPDMIA and INSULIN each reaching maximum importance (1.00), suggesting that this fold's data composition or decision boundary emphasized lipid and pharmacological factors over glycemic control alone. CKD showed moderate importance in Models 0 and 2 (0.34 and 0.21), while HPN and PAOD contributed minimally. The cross-model consistency of HbA1c reinforces its primacy as the most actionable counterfactual driver in DPN pre-diagnostic screening.}
% \label{fig:globalcf}
% \end{figure}

% \paragraph{HbA1c as the Dominant Counterfactual Driver}
HbA1c modification found to be the feature that was most changed across all models, (55 out of 58 for Model 0; 22 out of 22 for Model 1; and 13 out of 14 for Model 2; and 1 out of 3 for Model 3). In total, this is an overwhelming majority of 91 of the 97 total counterfactuals at a rate of 93.9\%. This means the glycemic index is deemed by the models as the most accessible path towards reclassification among all features identified as actionable.
This finding is clinically coherent and well-supported in the literature: chronic hyperglycemia—indexed by HbA1c—drives the oxidative and metabolic mechanisms underlying peripheral nerve damage in diabetic patients, and tighter glycemic control has been demonstrated to reduce DPN incidence and progression \cite{callaghan2012,tesfaye2010}. The counterfactual analysis, thus, not only identifies HbA1c as the most important model feature, but operationalizes this importance as an actionable clinical recommendation: patients flagged as high-risk by the model may be redirected toward glycemic optimization before, or in lieu of, a costly Nerve Conduction Study referral.

% Don-t include as it can be critiqued for lack of examples - Global counterfactual feature importance (Figure~\ref{fig:globalcf}) demonstrated a highly consistent pattern: HbA1c modification was necessary to reverse the DPN prediction in the overwhelming majority of counterfactuals across Models 0, 1, and 2 (importance scores: 0.95, 1.00, and 0.93, respectively). 

Model 3 exhibited a divergence from this pattern, with DSLPDMIA and INSULIN each attaining maximum counterfactual importance (3 out of 3) while HbA1c's accounted for 1 out of 3. Consideration has, of course, to be taken that there are only 3 counterfactuals for this model and all of them coming from a patient with an elevated level of HbA1c (12.4\%). This suggests that for patients in whom HbA1c is at a high level — dyslipidemia management and insulin therapy optimization may constitute the primary levers for risk reclassification.

% \paragraph{ Multifactorial Nature of Counterfactual Reversals}
Despite HbA1c's dominant role, the mean sparsity values across all four models ranged from 3.32 to 4.19 (Table~\ref{tab:localcf-model-level}), indicating that counterfactual reversals required changes to approximately three to four of the six available actionable features on average. This implies that the models do not rely on a single feature threshold to classify DPN risk, but rather on the combinatorial interaction of multiple comorbid conditions—a pattern consistent with the known multifactorial pathogenesis of DPN \cite{feldman2019}. 

CKD and dyslipidemia emerged as the second and third most frequently altered features (23 and 19 counterfactual instances, respectively (Table~\ref{tab:localcf-changed}), consistent with their established roles as peripheral neuropathy risk amplifiers through vascular and metabolic mechanisms \cite{jaiswal2017,callaghan2012}. 

In terms of distance from the reference instance, Model 0 has the highest mean among the models (L1 mean = 4.17; L2 = 3.29).  The higher distance values may be an effect of higher sparsity but also of the decision space of this model which allowed for higher classification margins.  Model 2 produced the most
parsimonious counterfactuals (L1 mean = 1.18; L2 mean = 0.82; sparsity mean = 3.32), suggesting that its decision boundary permits smaller, more targeted feature perturbations to achieve a prediction reversal.

The integration of counterfactual analysis into this pre-diagnostic screening framework addresses a critical barrier to ML adoption in clinical settings: interpretability. Black-box predictions, however accurate, are unlikely to be acted upon by clinicians without a justification that aligns with their clinical reasoning \cite{rudin2019}. By grounding model explanations in actionable, modifiable features and framing them as "what-if" scenarios consistent with clinical intervention pathways, the counterfactual approach offers a practical bridge between statistical prediction and clinical decision-making. Notably, the approach does not require clinicians to understand the internal mechanics of the CatBoost model; instead, it presents outputs in the form familiar to clinical reasoning: "If this patient's HbA1c and dyslipidemia were better controlled, the model would not flag them for NCS referral."

This has implications not only for individual patient management but for population-level screening program design, as the most frequently changed features across all counterfactuals—HbA1c, CKD status, and dyslipidemia—directly correspond to the targets of existing diabetes complication management guidelines \cite{ada2024}. Future work could explore whether counterfactual-informed clinical interventions in screened patients result in measurable changes in DPN incidence, providing prospective validation of the model's explanatory outputs.

\section{Conclusion}
%%%%%%%%%%%%%%%%%%%%%
This study demonstrates that a machine learning model — specifically CatBoost optimized through Bayesian hyperparameter search — can serve as a clinically viable pre-diagnostic screening tool for Diabetic Peripheral Neuropathy using routine clinical variables that do not require nerve conduction studies. Evaluated on 190 Filipino patients with Diabetes Mellitus across a repeated stratified cross-validation framework (10 repeats × 4 folds; 40 total runs), the final model achieved a mean sensitivity of 0.924 (95\% CI: 0.904–0.944) and an AUPRC of 0.885 (95\% CI: 0.868–0.902) — performance levels consistent with the clinical demands of a screening instrument whose primary function is to minimize missed DPN cases and rationalize referral to the more resource-intensive and geographically inaccessible NCS.

A distinguishing methodological contribution of this study is its systematic examination of DPN predictors across multiple complementary analytical lenses. Feature contribution was first assessed at the population level through a structured ablation study spanning twelve distinct feature set configurations, which revealed that neurological examination features produced the greatest performance degradation upon removal and that the combination of Sudoscan and neurological examination findings maintained performance close to, and in some instances exceeding, the full feature set. This was followed by CatBoost-native feature importance analysis and SHAP-based local interpretability, both of which converged consistently across all four cross-validation folds on decreased ankle reflex, decreased vibration sensation, Sudoscan-derived sudomotor parameters, and diabetes duration as the dominant predictors. Finally, counterfactual explanation analysis through the DiCE framework provided a fourth and clinically distinct analytical perspective — one focused not on what the model weighs most heavily, but on what can realistically be changed to alter a patient's predicted risk classification.

This multi-lens approach revealed an important tension inherent in conventional explainable AI methods as applied to clinical settings. While feature importance and SHAP analyses consistently identified neurological examination findings and Sudoscan measurements as among the most influential predictors — findings that are clinically plausible and biologically well-grounded — these features represent measurements of a patient's current neuropathic state rather than modifiable targets for intervention. A clinician cannot directly alter a patient's ankle reflex or electrochemical skin conductance through a clinical decision; these findings describe the disease, but do not point toward treatment. Counterfactual explanation analysis bridges this interpretability gap by reframing the question from "which features matter most?" to "which modifiable features, if changed, would alter the model's prediction?" In doing so, it shifts the clinical utility of the model from passive classification to active decision support.

Through this lens, HbA1c modification emerged as the dominant actionable pathway in 93.8\% of generated counterfactual instances across all four models, operationalizing glycemic optimization as the most accessible and clinically meaningful intervention for modifying model-predicted DPN risk. Dyslipidemia management and chronic kidney disease status emerged as secondary intervention targets, appearing in 19.6\% and 23.7\% of counterfactuals respectively — consistent with their established roles as peripheral neuropathy risk amplifiers through vascular and metabolic mechanisms. Notably, in patients with markedly elevated glycemic indices, insulin therapy initiation was additionally identified as a relevant therapeutic lever. These counterfactual outputs are particularly well-suited to the deployment context envisioned by this study: rather than simply flagging a patient for NCS referral, the system can simultaneously guide the clinician toward modifiable risk reduction strategies, potentially delaying or preventing neuropathic progression in borderline cases.

Several findings of this study are notable and merit further emphasis. The synergy between small-fiber sudomotor and large-fiber neurological biomarkers — demonstrated through both the feature set ablation results and the consistent co-prominence of these two feature groups across feature importance and SHAP analyses — suggests the model integrates complementary information spanning the full spectrum of peripheral nerve fiber subtypes affected in DPN. This multimodal integration is not commonly demonstrated in prior ML-based DPN studies, which have predominantly relied on routine demographic and glycemic variables. Furthermore, this study contributes to the geographic diversity of existing DPN-ML research by providing clinical evidence from the Philippine setting, where DPN affects 42–58\% of diabetic patients and where access to electrodiagnostic facilities remains uneven across regions.

Several limitations of this study must be explicitly acknowledged. The dataset was drawn from a single tertiary referral center — East Avenue Medical Center in Quezon City, Philippines — which may introduce institutional selection bias and limit the generalizability of findings to primary care or rural settings where patient phenotypes, comorbidity prevalence, and clinical resources may differ substantially. The relatively modest sample size of 190 patients, while addressed methodologically through repeated stratified cross-validation, constrains the statistical precision of performance estimates — particularly for specificity, which exhibited the widest confidence interval among all reported metrics. The modest class imbalance (approximately 69\% Confirmed DPN) further limits the precision of threshold-dependent metrics and the stability of counterfactual generation, with only 10 of 53 candidate instances successfully yielding counterfactuals within the defined computational timeout. While the absence of external validation is common in exploratory ML studies of this scale, it represents the most critical limitation with respect to clinical readiness: performance estimates derived entirely from internal cross-validation cannot substitute for prospective evaluation in an independent, geographically distinct cohort. Additionally, the counterfactual framework was restricted to a subset of modifiable features; features such as neurological examination findings and Sudoscan readings — despite their high predictive importance — could not be incorporated into actionable counterfactuals, as they reflect clinical measurements rather than interventional targets.

Future work should prioritize several directions. Prospective multi-center validation, ideally spanning both urban and rural Philippine clinical settings and extending to other Southeast Asian populations with comparable DPN prevalence profiles, is essential to establish the external validity and deployment readiness of the proposed screening tool. Expansion of the training dataset would improve model stability, enhance the precision of counterfactual generation, and enable more reliable evaluation of model performance in clinically important subgroups — including younger patients, women, and patients with near-normal glycemic profiles, who may present atypically and whose DPN risk may be underestimated by models trained on predominantly hyperglycemic cohorts. The relationship between chronic kidney disease and DPN classification warrants dedicated investigation: CKD appeared as the second most frequently modified feature in the counterfactual analysis, and its systematic appearance across multiple patients in a model for one data split (Model 0) suggests a potentially meaningful co-occurrence pattern between CKD status and confirmed DPN that the model learned from the training data. Whether this reflects a genuine biological interaction — whereby shared mechanisms of microvascular damage, uremic neuropathy, and metabolic dysregulation create compounding neuropathic risk in patients with concurrent CKD and diabetes — or a dataset-specific distributional artifact, cannot be resolved from the present study and merits targeted epidemiological and mechanistic investigation in larger, purpose-designed cohorts. Prospective investigation of whether counterfactual-guided clinical interventions — particularly glycemic optimization, dyslipidemia management, and nephroprotective strategies in patients flagged as high-risk — produce measurable reductions in DPN incidence or progression would provide direct prospective validation of the model's explanatory outputs. Integration of the screening framework into electronic medical record systems or mobile clinical decision support tools could facilitate real-world evaluation of its utility and clinician acceptability in point-of-care settings. Finally, future studies may explore whether the inclusion of additional modalities — such as plantar pressure measurements, optical coherence tomography of corneal nerve fibers, or continuous glucose monitoring-derived glycemic variability metrics — can further enhance screening sensitivity for the atypical DPN phenotypes that posed the greatest challenge to the present model.
In summary, this study demonstrates that an interpretable, counterfactual-augmented CatBoost model trained on accessible non-NCS clinical features can provide high-sensitivity DPN pre-diagnostic screening while simultaneously generating actionable, patient-specific clinical guidance. By examining DPN risk factors through the complementary lenses of feature set ablation, model-native feature importance, SHAP analysis, and counterfactual generation, the study offers a layered and clinically grounded account of how the model reaches its decisions and how those decisions can inform intervention. These findings establish a promising foundation for ML-assisted DPN triage in resource-constrained settings, and provide a structured framework for subsequent prospective validation and clinical deployment.
\backmatter

\bmhead{Acknowledgements}

The authors would like to thank Josiah Raziel S. Lluch of Mindanao State University -- Iligan Institute of Technology for his assistance with software implementation. This work is funded and monitored by the Department of Science and Technology (DOST) -- Philippine Council for Industry, Energy, and Emerging Technology Research and Development (PCIEERD) with Project No. 1213385.

\section*{Declarations}

%% NOTE: items below were populated only where the source manuscript stated
%% them explicitly (Funding, Ethics approval, Author contribution). The
%% remaining items are left as placeholders -- fill these in before
%% submission rather than assuming "Not applicable".

\begin{itemize}
\item Funding: This work is funded and monitored by the Department of Science and Technology (DOST) -- Philippine Council for Industry, Energy, and Emerging Technology Research and Development (PCIEERD) with Project No. 1213385.
\item Conflict of interest/Competing interests: [author to complete]
\item Ethics approval and consent to participate: This study was approved by the local Institutional Ethics Review Board and all subjects gave their written consent prior to participating.
\item Consent for publication: [author to complete]
\item Data availability: [author to complete]
\item Materials availability: [author to complete]
\item Code availability: [author to complete]
\item Author contribution: A.C.B. conducted the AI modeling and counterfactual generations and wrote the draft. E.M.M. and N.F. co-wrote portions of the draft. D.A.J., E.O.C., R.R.F., D.A.B., L.S.R., and M.T.O. designed and implemented the data collection process. N.F., L.D., and M.T.O. provided medical interpretations and insights. M.T.O. and C.P.M. provided research supervision. All authors reviewed the manuscript.
\end{itemize}

\begin{appendices}

\section{Counterfactuals Generated by Models 1--3}\label{secA1}

Figures~\ref{fig:localcf-m1-30}--\ref{fig:localcf-m3-97} show the counterfactual heatmaps for Models 1--3. They are provided without further explanation due to space limitations.

\begin{figure}[p]
\centering
\includegraphics[width=0.97\textwidth]{figures2/counterfactuals/030catboost_split1_local_cf.png}
\caption{Counterfactuals for Patient 30}
\label{fig:localcf-m1-30}
\end{figure}

\begin{figure}[p]
\centering
\includegraphics[width=0.97\textwidth]{figures2/counterfactuals/058catboost_split1_local_cf.png}
\caption{Counterfactuals for Patient 58}
\label{fig:localcf-m1-58}
\end{figure}

\begin{figure}[p]
\centering
\includegraphics[width=0.97\textwidth]{figures2/counterfactuals/074catboost_split1_local_cf.png}
\caption{Counterfactuals for Patient 74}
\label{fig:localcf-m1-74}
\end{figure}

\begin{figure}[p]
\centering
\includegraphics[width=0.97\textwidth]{figures2/counterfactuals/048catboost_split2_local_cf.png}
\caption{Counterfactuals for Patient 48}
\label{fig:localcf-m2-48}
\end{figure}

\begin{figure}[p]
\centering
\includegraphics[width=0.97\textwidth]{figures2/counterfactuals/055catboost_split2_local_cf.png}
\caption{Counterfactuals for Patient 55}
\label{fig:localcf-m2-55}
\end{figure}

\begin{figure}[p]
\centering
\includegraphics[width=0.97\textwidth]{figures2/counterfactuals/097catboost_split3_local_cf.png}
\caption{Counterfactuals for Patient 97}
\label{fig:localcf-m3-97}
\end{figure}

\end{appendices}

\bibliography{bibfile}

\end{document}
